% Options for packages loaded elsewhere
\PassOptionsToPackage{unicode}{hyperref}
\PassOptionsToPackage{hyphens}{url}
%
\documentclass[
]{article}
\usepackage{amsmath,amssymb}
\usepackage{iftex}
\ifPDFTeX
  \usepackage[T1]{fontenc}
  \usepackage[utf8]{inputenc}
  \usepackage{textcomp} % provide euro and other symbols
\else % if luatex or xetex
  \usepackage{unicode-math} % this also loads fontspec
  \defaultfontfeatures{Scale=MatchLowercase}
  \defaultfontfeatures[\rmfamily]{Ligatures=TeX,Scale=1}
\fi
\usepackage{lmodern}
\ifPDFTeX\else
  % xetex/luatex font selection
\fi
% Use upquote if available, for straight quotes in verbatim environments
\IfFileExists{upquote.sty}{\usepackage{upquote}}{}
\IfFileExists{microtype.sty}{% use microtype if available
  \usepackage[]{microtype}
  \UseMicrotypeSet[protrusion]{basicmath} % disable protrusion for tt fonts
}{}
\makeatletter
\@ifundefined{KOMAClassName}{% if non-KOMA class
  \IfFileExists{parskip.sty}{%
    \usepackage{parskip}
  }{% else
    \setlength{\parindent}{0pt}
    \setlength{\parskip}{6pt plus 2pt minus 1pt}}
}{% if KOMA class
  \KOMAoptions{parskip=half}}
\makeatother
\usepackage{xcolor}
\usepackage[margin=1in]{geometry}
\usepackage{longtable,booktabs,array}
\usepackage{calc} % for calculating minipage widths
% Correct order of tables after \paragraph or \subparagraph
\usepackage{etoolbox}
\makeatletter
\patchcmd\longtable{\par}{\if@noskipsec\mbox{}\fi\par}{}{}
\makeatother
% Allow footnotes in longtable head/foot
\IfFileExists{footnotehyper.sty}{\usepackage{footnotehyper}}{\usepackage{footnote}}
\makesavenoteenv{longtable}
\usepackage{graphicx}
\makeatletter
\newsavebox\pandoc@box
\newcommand*\pandocbounded[1]{% scales image to fit in text height/width
  \sbox\pandoc@box{#1}%
  \Gscale@div\@tempa{\textheight}{\dimexpr\ht\pandoc@box+\dp\pandoc@box\relax}%
  \Gscale@div\@tempb{\linewidth}{\wd\pandoc@box}%
  \ifdim\@tempb\p@<\@tempa\p@\let\@tempa\@tempb\fi% select the smaller of both
  \ifdim\@tempa\p@<\p@\scalebox{\@tempa}{\usebox\pandoc@box}%
  \else\usebox{\pandoc@box}%
  \fi%
}
% Set default figure placement to htbp
\def\fps@figure{htbp}
\makeatother
\setlength{\emergencystretch}{3em} % prevent overfull lines
\providecommand{\tightlist}{%
  \setlength{\itemsep}{0pt}\setlength{\parskip}{0pt}}
\setcounter{secnumdepth}{-\maxdimen} % remove section numbering
\usepackage{graphicx}
\usepackage{booktabs}
\usepackage{float}
\usepackage{titling}

% Logo's bovenaan
\pretitle{%
  \begin{center}
  \includegraphics[width=4cm]{figures/BINF.png}\hfill
  \includegraphics[width=6cm]{figures/PAIH.png}\\[2cm]
  \huge
}
\posttitle{\end{center}}
\usepackage{booktabs}
\usepackage{longtable}
\usepackage{array}
\usepackage{multirow}
\usepackage{wrapfig}
\usepackage{float}
\usepackage{colortbl}
\usepackage{pdflscape}
\usepackage{tabu}
\usepackage{threeparttable}
\usepackage{threeparttablex}
\usepackage[normalem]{ulem}
\usepackage{makecell}
\usepackage{xcolor}
\usepackage{bookmark}
\IfFileExists{xurl.sty}{\usepackage{xurl}}{} % add URL line breaks if available
\urlstyle{same}
\hypersetup{
  pdftitle={In-depth Analysis for Cathlab},
  hidelinks,
  pdfcreator={LaTeX via pandoc}}

\title{In-depth Analysis for Cathlab}
\usepackage{etoolbox}
\makeatletter
\providecommand{\subtitle}[1]{% add subtitle to \maketitle
  \apptocmd{\@title}{\par {\large #1 \par}}{}{}
}
\makeatother
\subtitle{Leveraging Data Analytics for Efficient OR Planning at ZOL}
\author{}
\date{\vspace{-2.5em}October 2025}

\begin{document}
\maketitle

\begin{center}
\vspace{2cm}

Niels MARTIN\\
niels.martin@uhasselt.be\\[0.5em]
Maxim RIEBUS\\
maxim.riebus@uhasselt.be\\[0.5em]
Haroon THARWAT\\
haroon.tharwat@uhasselt.be\\

\vspace{6cm}

\textbf{Hasselt University} \\
Faculty of Business Economics \\
Research Group Business Informatics \\
Process Analytics in Healthcare
\end{center}

\newpage
\tableofcontents
\newpage

\section{Introduction}\label{introduction}

Efficient management of operating rooms (ORs) requires an understanding
of how planned surgical activity aligns with real-world execution. This
document provides an in-depth, site-specific analysis for \textbf{Campus
Cathlab}, serving as an analytical extension of the previous exploratory
study. The goal is to systematically address the key research questions
formulated in the project description:

\textbf{Research questions}

\begin{itemize}
\tightlist
\item
  \textbf{Planned versus observed timing}

  \begin{itemize}
  \tightlist
  \item
    How well do the planned operating durations correspond to the actual
    durations?\\
  \item
    What is the typical deviation between the planned start time and the
    actual entry into the operating room?\\
  \item
    Which specialisms or procedure types show the largest differences?\\
  \item
    How do these deviations scale relative to the planned duration
    (proportional or percentage-wise)?
  \end{itemize}
\item
  \textbf{Adherence to working hours}

  \begin{itemize}
  \tightlist
  \item
    How often do procedures run beyond regular working hours (08:00 to
    16:30)?\\
  \item
    What proportion of OR activity occurs in the evening or at night?\\
  \item
    How do these patterns differ between campuses?
  \end{itemize}
\item
  \textbf{Room allocation}

  \begin{itemize}
  \tightlist
  \item
    How often are procedures performed in a different room than
    originally planned?\\
  \item
    Are these reallocations concentrated in specific rooms or
    specialisms?
  \end{itemize}
\item
  \textbf{Urgency and scheduling}

  \begin{itemize}
  \tightlist
  \item
    How are urgent and non-elective cases distributed across campuses
    and time periods?\\
  \item
    To what extent do urgent surgeries overlap with elective activity
    and cause scheduling disruptions?
  \end{itemize}
\end{itemize}

Each of the following sections addresses one of these themes in depth,
focusing on the local characteristics of the Cathlab campus.

\subsection{Unique values and data
richness}\label{unique-values-and-data-richness}

The dataset captures a high degree of heterogeneity in surgical
practice. A total of 7841 unique patient IDs are recorded, spread across
141 distinct procedure codes and 141 procedure names. Surgeon and
anesthesiologist coverage is broad, with 51 and 123 unique identifiers
respectively. These numbers reflect the diversity of clinical practice
at ZOL and underscore the importance of developing grouping strategies
to obtain actionable insights. Table \ref{tab:uniquevalues} provides an
overview of the number of unique values per variable.

\begin{table}[H]

\caption{\label{tab:unique values}Number of unique values per variable. \label{tab:uniquevalues}}
\centering
\begin{tabular}[t]{lr}
\toprule
Variable & UniqueValues\\
\midrule
PatientID & 7841\\
ProcedureCode & 141\\
ProcedureName & 141\\
AnesthesiologistID & 123\\
LengthOfStay & 117\\
\addlinespace
HeadSurgeonID & 51\\
PlannedOR & 12\\
Weekday & 7\\
ActualOR & 7\\
Year & 4\\
\addlinespace
AdmissionType & 3\\
\bottomrule
\end{tabular}
\end{table}

\subsection{Surgical procedure
durations}\label{surgical-procedure-durations}

Operating time is one of the central variables for planning in the
CathLab. Procedure durations exhibit substantial variation, both within
and across procedure types. Percentile analysis shows that 95 percent of
all procedures finish within 166 minutes, 98 percent within 204 minutes,
99 percent within 232 minutes, and 99.9 percent within 361.033 minutes.
Although most procedures are markedly shorter, the pronounced right-hand
tail of the distribution highlights the presence of outliers and a
limited number of highly complex cases that contribute
disproportionately to overall variability.

Figure \ref{fig:duration} visualises the distribution of procedure
durations, while Table \ref{tab:durationpercentiles} summarises the
corresponding percentile values. The combination of a concentrated bulk
and a long tail emphasises the operational challenge of balancing
predictable routine activity with rare but time-intensive interventions.

\begin{figure}[H]

{\centering \includegraphics[width=0.7\linewidth,]{In-Depth_Analysis_Cathlab_files/figure-latex/procedure_durations-1} 

}

\caption{Histogram of Surgery Duration \label{fig:duration}}\label{fig:procedure_durations}
\end{figure}

\newpage

\begin{longtable}[]{@{}lr@{}}
\caption{Percentiles of surgical procedure duration (in minutes).
\label{tab:durationpercentiles}}\tabularnewline
\toprule\noalign{}
Percentile & Duration (minutes) \\
\midrule\noalign{}
\endfirsthead
\toprule\noalign{}
Percentile & Duration (minutes) \\
\midrule\noalign{}
\endhead
\bottomrule\noalign{}
\endlastfoot
95\% & 166.000 \\
98\% & 204.000 \\
99\% & 232.000 \\
99.9\% & 361.033 \\
\end{longtable}

Examining the most frequent procedures reveals substantial variation in
operating times. ABLATIE VKF, the most common procedure, has an average
duration of 98.1 minutes and a median of 99 minutes, while DILATATIE
FEMORALIS averages 80.0 minutes with a median of 74.5 minutes. EFO
DIAGNOSTISCH is considerably shorter, with a mean of 56.1 minutes and a
median of 53 minutes. At the other end of the spectrum, RX
interventionele neuro procedures show markedly longer durations,
averaging 150.9 minutes with a median of 143 minutes. This contrast
highlights the different resource demands associated with routine,
standardised interventions versus more complex and time-intensive cases.

The coexistence of high-volume short procedures and less frequent but
substantially longer interventions illustrates the operational challenge
of designing a schedule that remains both efficient and flexible. Table
\ref{tab:top10dur} summarises the ten most common procedure types and
their corresponding duration characteristics.

\begin{longtable}[]{@{}lrrr@{}}
\caption{Top 10 most frequent procedure types with mean and median
duration \label{tab:top10dur}}\tabularnewline
\toprule\noalign{}
ProcedureName & Count & Mean & Median \\
\midrule\noalign{}
\endfirsthead
\toprule\noalign{}
ProcedureName & Count & Mean & Median \\
\midrule\noalign{}
\endhead
\bottomrule\noalign{}
\endlastfoot
ABLATIE VKF & 2251 & 98.1 & 99.0 \\
DILATATIE FEMORALIS & 818 & 80.0 & 74.5 \\
EFO DIAGNOSTISCH & 543 & 56.1 & 53.0 \\
DILATATIE ILIACA & 466 & 70.0 & 65.0 \\
RX interventionele neuro & 449 & 150.9 & 143.0 \\
ABLATIE FLUTTER & 357 & 65.8 & 63.0 \\
ABLATIE AV NODALE REENTRY TACHYCARDIE & 350 & 71.3 & 68.0 \\
DILATATIE TIBIALIS & 273 & 98.3 & 93.0 \\
ABLATIE BUNDEL HIS & 237 & 46.6 & 43.0 \\
Percutane aortaklep vervanging & 237 & 122.8 & 118.0 \\
\end{longtable}

\subsection{Admission type and
urgency}\label{admission-type-and-urgency}

The majority of procedures in the dataset occur in a planned setting.
Ambulatory admissions represent 88.6\% of all cases, while inpatient
hospitalizations account for 11.3\%. Emergency admissions are rare,
comprising only 0.1\% of recorded activity. The urgency classification
mirrors this distribution: 97.5\% of procedures are elective, and 2.5\%
are non-elective. Although the latter group is relatively small, it
remains operationally relevant, as urgent cases can disrupt elective
schedules, reduce planning stability, and increase the likelihood of
cancellations or overtime. Table \ref{tab:admissiontype} and Table
\ref{tab:urgencytype} summarize both dimensions of case mix.

\begin{longtable}[]{@{}lrr@{}}
\caption{Frequency table for Admission Type
\label{tab:admissiontype}}\tabularnewline
\toprule\noalign{}
AdmissionType & Count & Percentage \\
\midrule\noalign{}
\endfirsthead
\toprule\noalign{}
AdmissionType & Count & Percentage \\
\midrule\noalign{}
\endhead
\bottomrule\noalign{}
\endlastfoot
HOS & 8222 & 88.6 \\
DAG & 1053 & 11.3 \\
SPOED & 7 & 0.1 \\
\end{longtable}

\begin{longtable}[]{@{}lrr@{}}
\caption{Frequency table for Urgency Classification
\label{tab:urgencytype}}\tabularnewline
\toprule\noalign{}
UrgencyType & Count & Percentage \\
\midrule\noalign{}
\endfirsthead
\toprule\noalign{}
UrgencyType & Count & Percentage \\
\midrule\noalign{}
\endhead
\bottomrule\noalign{}
\endlastfoot
electief & 9050 & 97.5 \\
niet-electief & 232 & 2.5 \\
\end{longtable}

\subsection{Temporal distribution of
activity}\label{temporal-distribution-of-activity}

Surgical activity is distributed consistently across weekdays, with
Wednesday representing the largest share at 24.9 percent of all
procedures, followed by Thursday at 21.9 percent, Tuesday at 19.5
percent, and Friday at 17.2 percent. Monday accounts for 15.1 percent of
weekly activity. Weekend activity is limited but present, with 0.8
percent of surgeries performed on Sunday and 0.6 percent on Saturday.
Annual volumes show steady growth, increasing from 2571 procedures in
2022 to 2858 in 2024. In the partial dataset for 2025, 1199 cases
represent 12.9 percent of the total activity. Together, Table
\ref{tab:weekday} and Table \ref{tab:year} summarize the temporal
distribution of procedures across weekdays and years.

\begin{longtable}[]{@{}
  >{\centering\arraybackslash}p{(\linewidth - 4\tabcolsep) * \real{0.1389}}
  >{\centering\arraybackslash}p{(\linewidth - 4\tabcolsep) * \real{0.1111}}
  >{\centering\arraybackslash}p{(\linewidth - 4\tabcolsep) * \real{0.1806}}@{}}
\caption{Frequency table for Day of the Week (sorted by descending
count) \label{tab:weekday}}\tabularnewline
\toprule\noalign{}
\begin{minipage}[b]{\linewidth}\centering
Weekday
\end{minipage} & \begin{minipage}[b]{\linewidth}\centering
Count
\end{minipage} & \begin{minipage}[b]{\linewidth}\centering
Percentage
\end{minipage} \\
\midrule\noalign{}
\endfirsthead
\toprule\noalign{}
\begin{minipage}[b]{\linewidth}\centering
Weekday
\end{minipage} & \begin{minipage}[b]{\linewidth}\centering
Count
\end{minipage} & \begin{minipage}[b]{\linewidth}\centering
Percentage
\end{minipage} \\
\midrule\noalign{}
\endhead
\bottomrule\noalign{}
\endlastfoot
Wo & 2308 & 24.9 \\
Do & 2034 & 21.9 \\
Di & 1812 & 19.5 \\
Vr & 1598 & 17.2 \\
Ma & 1398 & 15.1 \\
Zo & 73 & 0.8 \\
Za & 59 & 0.6 \\
\end{longtable}

\begin{longtable}[]{@{}
  >{\centering\arraybackslash}p{(\linewidth - 4\tabcolsep) * \real{0.0972}}
  >{\centering\arraybackslash}p{(\linewidth - 4\tabcolsep) * \real{0.1111}}
  >{\centering\arraybackslash}p{(\linewidth - 4\tabcolsep) * \real{0.1806}}@{}}
\caption{Frequency table for Year \label{tab:year}}\tabularnewline
\toprule\noalign{}
\begin{minipage}[b]{\linewidth}\centering
Year
\end{minipage} & \begin{minipage}[b]{\linewidth}\centering
Count
\end{minipage} & \begin{minipage}[b]{\linewidth}\centering
Percentage
\end{minipage} \\
\midrule\noalign{}
\endfirsthead
\toprule\noalign{}
\begin{minipage}[b]{\linewidth}\centering
Year
\end{minipage} & \begin{minipage}[b]{\linewidth}\centering
Count
\end{minipage} & \begin{minipage}[b]{\linewidth}\centering
Percentage
\end{minipage} \\
\midrule\noalign{}
\endhead
\bottomrule\noalign{}
\endlastfoot
2022 & 2571 & 27.7 \\
2023 & 2654 & 28.6 \\
2024 & 2858 & 30.8 \\
2025 & 1199 & 12.9 \\
\end{longtable}

In summary, procedure durations exhibit substantial variability, with a
long-tailed distribution driven by more complex interventions. Most
procedures are elective, yet the smaller share of non-elective cases
remains operationally significant because of its potential to disrupt
schedule stability. Surgical activity is spread evenly across weekdays
and has increased steadily over recent years. Together, these patterns
establish a clear baseline for the dataset and provide essential context
for subsequent analyses of start and end time deviations, utilization,
and predictive modelling in the following sections of this report.

\newpage

\section{How accurately does planning reflect reality in the
OR?}\label{how-accurately-does-planning-reflect-reality-in-the-or}

\subsection{Are planned and observed durations
aligned?}\label{are-planned-and-observed-durations-aligned}

The comparison between planned and observed surgery durations shows
considerable dispersion around the one-to-one line. Many procedures fall
close to the diagonal, suggesting that their observed durations align
reasonably well with planning estimates. However, a substantial share of
cases lies noticeably above or below this line, confirming that both
underestimation and overestimation are common in practice. Deviations
are more often in the direction of longer-than-planned surgeries,
reflected by the dense cluster of points below the red dashed line. The
spread widens for longer procedures, illustrating the increased
uncertainty associated with more complex interventions. Figure
\ref{fig:planned_vs_observed} visualizes this relationship.

\begin{figure}[h]

{\centering \includegraphics[width=0.7\linewidth,]{In-Depth_Analysis_Cathlab_files/figure-latex/planned vs observed-1} 

}

\caption{Planned vs. Observed Surgery Duration (Cathlab) \label{fig:planned_vs_observed}}\label{fig:planned vs observed}
\end{figure}

On average, surgeries at Cathlab show substantial variation between
planned and actual durations. For procedures that run longer than
scheduled, the mean deviation is 26.9 minutes with a median of 18
minutes. Cases that finish earlier deviate by an average of 25.7 minutes
and a median of 15 minutes. The corresponding standard deviations, 29.4
minutes for longer cases and 82.1 minutes for shorter ones, indicate
that while many deviations are moderate, a small subset of procedures
deviates by a wide margin. These outliers extend the overall range of
planning differences and illustrate the inherent unpredictability of
certain interventions.

Almost half of all surgeries (43.0 percent) exceed their planned
duration, with an interquartile range of 27 minutes, whereas 57.0
percent finish ahead of schedule with an interquartile range of 19
minutes. Deviations therefore occur in both directions, but the spread
is noticeably larger among earlier finishing cases. This variability
suggests that although planning accuracy is acceptable when viewed in
aggregate, predicting individual case durations remains challenging.
Enhancing the stability of duration estimates, particularly for
procedures that frequently deviate, could help reduce day-to-day
variability in the schedule and limit downstream delays across the
operating list. Table \ref{tab:dur_diff_direction} provides a detailed
summary of these patterns.

\begin{longtable}[]{@{}
  >{\centering\arraybackslash}p{(\linewidth - 20\tabcolsep) * \real{0.2347}}
  >{\centering\arraybackslash}p{(\linewidth - 20\tabcolsep) * \real{0.0714}}
  >{\centering\arraybackslash}p{(\linewidth - 20\tabcolsep) * \real{0.0918}}
  >{\centering\arraybackslash}p{(\linewidth - 20\tabcolsep) * \real{0.0714}}
  >{\centering\arraybackslash}p{(\linewidth - 20\tabcolsep) * \real{0.0612}}
  >{\centering\arraybackslash}p{(\linewidth - 20\tabcolsep) * \real{0.0612}}
  >{\centering\arraybackslash}p{(\linewidth - 20\tabcolsep) * \real{0.0714}}
  >{\centering\arraybackslash}p{(\linewidth - 20\tabcolsep) * \real{0.0714}}
  >{\centering\arraybackslash}p{(\linewidth - 20\tabcolsep) * \real{0.0816}}
  >{\centering\arraybackslash}p{(\linewidth - 20\tabcolsep) * \real{0.0816}}
  >{\centering\arraybackslash}p{(\linewidth - 20\tabcolsep) * \real{0.1020}}@{}}
\caption{Difference between actual and planned duration (minutes),
separated by direction \label{tab:dur_diff_direction}}\tabularnewline
\toprule\noalign{}
\begin{minipage}[b]{\linewidth}\centering
direction
\end{minipage} & \begin{minipage}[b]{\linewidth}\centering
Mean
\end{minipage} & \begin{minipage}[b]{\linewidth}\centering
Median
\end{minipage} & \begin{minipage}[b]{\linewidth}\centering
SD
\end{minipage} & \begin{minipage}[b]{\linewidth}\centering
IQR
\end{minipage} & \begin{minipage}[b]{\linewidth}\centering
Min
\end{minipage} & \begin{minipage}[b]{\linewidth}\centering
Max
\end{minipage} & \begin{minipage}[b]{\linewidth}\centering
P95
\end{minipage} & \begin{minipage}[b]{\linewidth}\centering
P99
\end{minipage} & \begin{minipage}[b]{\linewidth}\centering
P99\_9
\end{minipage} & \begin{minipage}[b]{\linewidth}\centering
Percent
\end{minipage} \\
\midrule\noalign{}
\endfirsthead
\toprule\noalign{}
\begin{minipage}[b]{\linewidth}\centering
direction
\end{minipage} & \begin{minipage}[b]{\linewidth}\centering
Mean
\end{minipage} & \begin{minipage}[b]{\linewidth}\centering
Median
\end{minipage} & \begin{minipage}[b]{\linewidth}\centering
SD
\end{minipage} & \begin{minipage}[b]{\linewidth}\centering
IQR
\end{minipage} & \begin{minipage}[b]{\linewidth}\centering
Min
\end{minipage} & \begin{minipage}[b]{\linewidth}\centering
Max
\end{minipage} & \begin{minipage}[b]{\linewidth}\centering
P95
\end{minipage} & \begin{minipage}[b]{\linewidth}\centering
P99
\end{minipage} & \begin{minipage}[b]{\linewidth}\centering
P99\_9
\end{minipage} & \begin{minipage}[b]{\linewidth}\centering
Percent
\end{minipage} \\
\midrule\noalign{}
\endhead
\bottomrule\noalign{}
\endlastfoot
Longer than planned & 26.9 & 18 & 29.4 & 27 & 1 & 350 & 81.3 & 140 & 261
& 43.0\% \\
Shorter than planned & 25.7 & 15 & 82.1 & 19 & 0 & 1363 & 50 & 118.4 &
1184 & 57.0\% \\
\end{longtable}

Expressing deviations as a percentage of the planned duration provides a
standardized perspective that allows short and long procedures to be
evaluated on the same scale. Most cases fall within approximately ±40
percent of their planned time, indicating a broadly acceptable level of
scheduling accuracy for day-to-day operations. The histogram peaks just
below zero, showing that procedures tend to finish slightly earlier than
scheduled on average. The distribution is roughly symmetrical near the
center, though the right tail stretches farther than the left,
reflecting that while major overruns are relatively infrequent, they can
be substantial when they do occur. This pattern suggests that planning
accuracy is reasonable for the majority of cases, but a smaller subset
of surgeries contributes disproportionately to the overall spread in
relative deviations. Figure \ref{fig:relative_dev} displays this
distribution.

\begin{figure}[h]

{\centering \includegraphics[width=0.7\linewidth,]{In-Depth_Analysis_Cathlab_files/figure-latex/relative_deviation_hist-1} 

}

\caption{Distribution of relative deviation between planned and observed duration \label{fig:relative_dev}}\label{fig:relative_deviation_hist}
\end{figure}

Duration deviations differ markedly between admission types. Emergency,
SPOED, cases display the highest unpredictability. Their mean relative
deviation is minus 11.3 percent, with a median of minus 8 percent,
indicating that these procedures usually finish earlier than planned.
The large spread in the distribution shows that individual cases can
deviate substantially from their planned duration, which is
characteristic of acute, unplanned interventions where both clinical
urgency and case complexity vary widely.

Inpatient, HOS, procedures show more moderate deviation patterns. The
mean relative deviation of 3.7 percent and the median of minus 3.5
percent suggest that most surgeries fall reasonably close to plan, with
a slight tendency to finish earlier. However, the variability remains
considerable, as reflected by the wide distribution of deviations. This
indicates that while the central tendency aligns with planning
expectations, a notable subset of cases still diverges meaningfully from
their scheduled durations.

Ambulatory, DAG, procedures are the most predictable group. Their mean
deviation is minus 2.4 percent and the median minus 8.8 percent, showing
that short stay procedures are generally completed slightly ahead of
schedule. The relatively large percentage deviations observed here arise
mainly because even small absolute differences represent a sizable
proportion of the short planned durations typical of this category.
Overall, these cases exhibit the tightest duration patterns, reflecting
the standardized and well structured nature of ambulatory workflows.

Table \ref{tab:rel_dev_admission} summarizes these patterns across
admission types.

\begin{longtable}[]{@{}
  >{\centering\arraybackslash}p{(\linewidth - 10\tabcolsep) * \real{0.1553}}
  >{\centering\arraybackslash}p{(\linewidth - 10\tabcolsep) * \real{0.0680}}
  >{\centering\arraybackslash}p{(\linewidth - 10\tabcolsep) * \real{0.2913}}
  >{\centering\arraybackslash}p{(\linewidth - 10\tabcolsep) * \real{0.3107}}
  >{\centering\arraybackslash}p{(\linewidth - 10\tabcolsep) * \real{0.0874}}
  >{\centering\arraybackslash}p{(\linewidth - 10\tabcolsep) * \real{0.0874}}@{}}
\caption{Relative duration deviation and coefficient of variation by
admission type (Cathlab) \label{tab:rel_dev_admission}}\tabularnewline
\toprule\noalign{}
\begin{minipage}[b]{\linewidth}\centering
AdmissionType
\end{minipage} & \begin{minipage}[b]{\linewidth}\centering
n
\end{minipage} & \begin{minipage}[b]{\linewidth}\centering
Mean relative deviation (\%)
\end{minipage} & \begin{minipage}[b]{\linewidth}\centering
Median relative deviation (\%)
\end{minipage} & \begin{minipage}[b]{\linewidth}\centering
SD (\%)
\end{minipage} & \begin{minipage}[b]{\linewidth}\centering
CV (\%)
\end{minipage} \\
\midrule\noalign{}
\endfirsthead
\toprule\noalign{}
\begin{minipage}[b]{\linewidth}\centering
AdmissionType
\end{minipage} & \begin{minipage}[b]{\linewidth}\centering
n
\end{minipage} & \begin{minipage}[b]{\linewidth}\centering
Mean relative deviation (\%)
\end{minipage} & \begin{minipage}[b]{\linewidth}\centering
Median relative deviation (\%)
\end{minipage} & \begin{minipage}[b]{\linewidth}\centering
SD (\%)
\end{minipage} & \begin{minipage}[b]{\linewidth}\centering
CV (\%)
\end{minipage} \\
\midrule\noalign{}
\endhead
\bottomrule\noalign{}
\endlastfoot
HOS & 8222 & 3.7 & -3.5 & 40.1 & 10.9 \\
DAG & 1053 & -2.4 & -8.8 & 39 & -16.1 \\
SPOED & 7 & -11.3 & -8 & 32.2 & -2.9 \\
\end{longtable}

A similar pattern appears when comparing elective and non elective
surgeries. Non elective procedures show the largest deviations, with an
average relative deviation of 16.7 percent and a median of 4.4 percent.
These positive values indicate that urgent cases often run longer than
planned. Their relatively low coefficient of variation reflects the fact
that deviations are consistently large across cases rather than
dominated by a few extremes, underscoring the unpredictable nature of
non elective activity.

Elective operations, by contrast, show much smaller deviations: the mean
relative deviation is 2.6 percent, and the median is minus 4.3 percent,
indicating that many elective procedures finish slightly earlier than
expected. The higher coefficient of variation in this group stems from a
small number of substantial outliers that widen the distribution, even
though most elective surgeries remain close to their planned durations.

Together, these findings confirm that elective cases are generally well
estimated, while non elective surgeries introduce the greatest
variability and therefore pose the biggest challenge to maintaining a
stable operating schedule. Table \ref{tab:rel_dev_urgency} summarizes
these results.

\begin{longtable}[]{@{}
  >{\centering\arraybackslash}p{(\linewidth - 10\tabcolsep) * \real{0.1553}}
  >{\centering\arraybackslash}p{(\linewidth - 10\tabcolsep) * \real{0.0680}}
  >{\centering\arraybackslash}p{(\linewidth - 10\tabcolsep) * \real{0.2913}}
  >{\centering\arraybackslash}p{(\linewidth - 10\tabcolsep) * \real{0.3107}}
  >{\centering\arraybackslash}p{(\linewidth - 10\tabcolsep) * \real{0.0874}}
  >{\centering\arraybackslash}p{(\linewidth - 10\tabcolsep) * \real{0.0874}}@{}}
\caption{Relative duration deviation and coefficient of variation by
urgency type (Cathlab) \label{tab:rel_dev_urgency}}\tabularnewline
\toprule\noalign{}
\begin{minipage}[b]{\linewidth}\centering
UrgencyType
\end{minipage} & \begin{minipage}[b]{\linewidth}\centering
n
\end{minipage} & \begin{minipage}[b]{\linewidth}\centering
Mean relative deviation (\%)
\end{minipage} & \begin{minipage}[b]{\linewidth}\centering
Median relative deviation (\%)
\end{minipage} & \begin{minipage}[b]{\linewidth}\centering
SD (\%)
\end{minipage} & \begin{minipage}[b]{\linewidth}\centering
CV (\%)
\end{minipage} \\
\midrule\noalign{}
\endfirsthead
\toprule\noalign{}
\begin{minipage}[b]{\linewidth}\centering
UrgencyType
\end{minipage} & \begin{minipage}[b]{\linewidth}\centering
n
\end{minipage} & \begin{minipage}[b]{\linewidth}\centering
Mean relative deviation (\%)
\end{minipage} & \begin{minipage}[b]{\linewidth}\centering
Median relative deviation (\%)
\end{minipage} & \begin{minipage}[b]{\linewidth}\centering
SD (\%)
\end{minipage} & \begin{minipage}[b]{\linewidth}\centering
CV (\%)
\end{minipage} \\
\midrule\noalign{}
\endhead
\bottomrule\noalign{}
\endlastfoot
niet-electief & 232 & 16.7 & 4.4 & 60.9 & 3.7 \\
electief & 9050 & 2.6 & -4.3 & 39.2 & 14.9 \\
\end{longtable}

The mean relative deviation between planned and observed durations
remains broadly stable across the observed years. In 2022, the average
deviation was 6.7 percent, decreasing to 4 percent in 2023. A more
pronounced decline occurred in 2024, when the mean deviation reached
minus 0.8 percent, while the partial data for 2025 show a modest
increase to 1.9 percent.

These year to year fluctuations are small and likely reflect normal
operational variability rather than systematic changes in planning
accuracy. Overall, the consistency of the yearly averages indicates that
the scheduling process at Cathlab is well calibrated at the aggregate
level. Despite this stability, individual procedures continue to exhibit
substantial variability, which naturally limits day to day precision.
Figure \ref{fig:reldev_year} illustrates these trends.

\begin{figure}[H]

{\centering \includegraphics[width=0.7\linewidth,]{In-Depth_Analysis_Cathlab_files/figure-latex/rel_duration_dev_year-1} 

}

\caption{Relative duration deviation by year \label{fig:reldev_year}}\label{fig:rel_duration_dev_year}
\end{figure}

The relative deviation between planned and observed durations remains
stable across the regular workweek, with mean deviations ranging from
minus 1.2 percent to 3.9 percent. This consistency suggests that
planning accuracy does not vary meaningfully from Monday to Friday.
Minor fluctuations, such as the slightly higher averages on Tuesday and
Wednesday, likely reflect differences in case mix rather than structural
scheduling inefficiencies.

A pronounced shift appears during the weekend. On Saturdays and Sundays,
mean deviations increase steeply to 16 percent and 23.8 percent, far
exceeding weekday levels. These larger discrepancies stem from the
smaller and more unpredictable volume of urgent cases. Weekend activity
is predominantly non elective, requiring flexible resource allocation
and exhibiting less predictable procedure durations. The resulting
contrast between weekdays and weekends underscores the stability of
elective scheduling during standard operating days and the inherent
variability associated with emergency work at the end of the week.
Figure \ref{fig:reldev_weekday} summarises these patterns.

\begin{figure}[H]

{\centering \includegraphics[width=0.7\linewidth,]{In-Depth_Analysis_Cathlab_files/figure-latex/rel_dev_weekday-1} 

}

\caption{Relative deviation by weekday \label{fig:reldev_weekday}}\label{fig:rel_dev_weekday}
\end{figure}

Variation in duration deviation differs considerably across operating
rooms, as shown in Table \ref{tab:ordeviation_or}. The largest
deviations occur in KO03, which stands out with a mean absolute
deviation of 112.4 minutes and 67.6 percent of cases running longer than
planned. KO06 and KO07 also exhibit substantial deviations, 48.3 and
30.9 minutes on average, reflecting the higher uncertainty and
procedural complexity characteristic of these rooms.

By contrast, rooms such as KO01 and KO02 show much smaller deviations,
with mean absolute values of 26.1 and 20 minutes respectively, and a
more balanced distribution between shorter and longer cases. These lower
deviations suggest that the procedures performed in these rooms are more
standardized or predictable, leading to greater planning accuracy.

Overall, the differences across rooms highlight how case mix and
procedural complexity strongly shape planning precision. Rooms that
handle diverse or technically demanding interventions show wider
deviation patterns, while those with more routine activity display
tighter and more predictable scheduling behaviour.

\begin{longtable}[]{@{}
  >{\centering\arraybackslash}p{(\linewidth - 12\tabcolsep) * \real{0.0924}}
  >{\centering\arraybackslash}p{(\linewidth - 12\tabcolsep) * \real{0.0588}}
  >{\centering\arraybackslash}p{(\linewidth - 12\tabcolsep) * \real{0.1261}}
  >{\centering\arraybackslash}p{(\linewidth - 12\tabcolsep) * \real{0.1092}}
  >{\centering\arraybackslash}p{(\linewidth - 12\tabcolsep) * \real{0.2437}}
  >{\centering\arraybackslash}p{(\linewidth - 12\tabcolsep) * \real{0.1176}}
  >{\centering\arraybackslash}p{(\linewidth - 12\tabcolsep) * \real{0.2521}}@{}}
\caption{Duration deviations per operating room, ordered by mean
absolute deviation \label{tab:ordeviation_or}}\tabularnewline
\toprule\noalign{}
\begin{minipage}[b]{\linewidth}\centering
ActualOR
\end{minipage} & \begin{minipage}[b]{\linewidth}\centering
n
\end{minipage} & \begin{minipage}[b]{\linewidth}\centering
Mean abs dev
\end{minipage} & \begin{minipage}[b]{\linewidth}\centering
Longer (\%)
\end{minipage} & \begin{minipage}[b]{\linewidth}\centering
Mean relative dev (longer)
\end{minipage} & \begin{minipage}[b]{\linewidth}\centering
Shorter (\%)
\end{minipage} & \begin{minipage}[b]{\linewidth}\centering
Mean relative dev (shorter)
\end{minipage} \\
\midrule\noalign{}
\endfirsthead
\toprule\noalign{}
\begin{minipage}[b]{\linewidth}\centering
ActualOR
\end{minipage} & \begin{minipage}[b]{\linewidth}\centering
n
\end{minipage} & \begin{minipage}[b]{\linewidth}\centering
Mean abs dev
\end{minipage} & \begin{minipage}[b]{\linewidth}\centering
Longer (\%)
\end{minipage} & \begin{minipage}[b]{\linewidth}\centering
Mean relative dev (longer)
\end{minipage} & \begin{minipage}[b]{\linewidth}\centering
Shorter (\%)
\end{minipage} & \begin{minipage}[b]{\linewidth}\centering
Mean relative dev (shorter)
\end{minipage} \\
\midrule\noalign{}
\endhead
\bottomrule\noalign{}
\endlastfoot
KO03 & 34 & 112.4 & 67.6 & 117.8 & 32.4 & -39 \\
KO06 & 961 & 48.3 & 59.6 & 50.8 & 39.3 & -27.4 \\
KO07 & 914 & 30.9 & 47.8 & 38.5 & 50.8 & -19.9 \\
KO05 & 2665 & 26.1 & 42.5 & 35.4 & 55.8 & -22.3 \\
KO04 & 10 & 23.1 & 60 & 112.7 & 40 & -31.2 \\
KO01 & 1846 & 20.5 & 36.3 & 26.2 & 61.6 & -19.2 \\
KO02 & 2852 & 20 & 40.5 & 25.3 & 57.5 & -20.4 \\
\end{longtable}

Among the most extreme deviations in surgery duration, longer and
shorter cases occur in nearly equal proportions. Of the top one percent
of deviations, 48.9 percent correspond to procedures that last longer
than planned, while 51.1 percent finish earlier than expected. This
close balance contrasts with earlier observations in a different
dataset, where overruns were more dominant.

Although both directions appear among the extreme cases, their
operational implications differ. Early finishes typically create limited
idle time, whereas overruns can delay subsequent cases and increase the
likelihood of overtime. The near symmetry between longer and shorter
cases in the extreme deviations suggests that both large
underestimations and large overestimations meaningfully contribute to
schedule variability. This reinforces the need to monitor procedures
prone to substantial deviations in either direction. Table
\ref{tab:extreme_dir} summarizes these results.

\begin{longtable}[]{@{}
  >{\centering\arraybackslash}p{(\linewidth - 4\tabcolsep) * \real{0.3194}}
  >{\centering\arraybackslash}p{(\linewidth - 4\tabcolsep) * \real{0.0694}}
  >{\centering\arraybackslash}p{(\linewidth - 4\tabcolsep) * \real{0.1806}}@{}}
\caption{Direction of deviation among the 1\% most extreme cases
(Cathlab) \label{tab:extreme_dir}}\tabularnewline
\toprule\noalign{}
\begin{minipage}[b]{\linewidth}\centering
Direction
\end{minipage} & \begin{minipage}[b]{\linewidth}\centering
n
\end{minipage} & \begin{minipage}[b]{\linewidth}\centering
Percentage
\end{minipage} \\
\midrule\noalign{}
\endfirsthead
\toprule\noalign{}
\begin{minipage}[b]{\linewidth}\centering
Direction
\end{minipage} & \begin{minipage}[b]{\linewidth}\centering
n
\end{minipage} & \begin{minipage}[b]{\linewidth}\centering
Percentage
\end{minipage} \\
\midrule\noalign{}
\endhead
\bottomrule\noalign{}
\endlastfoot
Longer than planned & 46 & 48.9\% \\
Shorter than planned & 48 & 51.1\% \\
\end{longtable}

\subsection{A closer look on the coefficient of
variation}\label{a-closer-look-on-the-coefficient-of-variation}

The monthly coefficient of variation for planning deviation fluctuates
clearly over time, with pronounced spikes and drops, but no sustained
long-term upward trend. As Figure \ref{fig:cv_planning} shows, early
2022 begins with relatively low values near 1.0 before increasing toward
mid-2022, where the coefficient briefly rises above 1.5 and reaches its
first major peak around 3.0. After this spike, the metric returns to
levels between roughly 1.2 and 1.6 toward the end of 2022.

Variability intensifies in 2023. Several sharp increases appear, with
multiple months reaching or exceeding 2.5 and occasional peaks above
3.0. These are interspersed with steep declines back toward values near
1.0, producing a sawtooth pattern rather than a stable trend. The
alternation of short high-variability episodes followed by abrupt dips
suggests that month-to-month case mix or scheduling conditions play a
major role in driving dispersion.

In 2024 the pattern continues, with some of the highest peaks in the
entire series---several months exceed 3.0, followed again by sudden
returns toward values near 1.0. The amplitude of these swings indicates
that some periods experience markedly unpredictable deviations, while
others revert to relatively stable planning performance.

The first months of 2025 remain highly variable as well: values
oscillate from just above 1.0 to more than 3.0 within short intervals.
This reinforces the absence of a stabilising trend and shows that
planning deviation variability continues to fluctuate strongly.

Overall, the evolution of the coefficient is dominated by repeated high
peaks and abrupt drops. The metric does not drift upward or downward
structurally but cycles between periods of elevated unpredictability and
periods of more controlled deviation. The recurring spikes above 3.0
indicate months with substantially more erratic scheduling accuracy,
while the frequent returns toward values near 1.0 show that these
episodes are temporary rather than systematic.

\begin{figure}[H]

{\centering \includegraphics[width=0.7\linewidth,]{In-Depth_Analysis_Cathlab_files/figure-latex/cv_monthly-1} 

}

\caption{Monthly evolution of coefficient of variation (Planning deviation) \label{fig:cv_planning}}\label{fig:cv_monthly}
\end{figure}

\subsubsection{Variability by planned duration
group}\label{variability-by-planned-duration-group}

Variation in surgical timing can be assessed from two complementary
angles, each capturing a different source of performance heterogeneity.
The first is the coefficient of variation of observed duration, defined
as the standard deviation divided by the mean of the recorded surgical
times. In a dataset containing both very short and very long operations,
this metric reflects the structural spread of procedure durations across
the case mix rather than instability within individual surgeries. High
values therefore stem primarily from the diversity of procedure types,
not from erratic execution.

The second perspective is the coefficient of variation of the planning
deviation. This measure takes the standard deviation of the absolute
difference between planned and observed duration and divides it by the
mean of that deviation. It expresses how consistently planning accuracy
is achieved. Unlike the CV of observed duration, which tracks intrinsic
variation in surgical work, the CV of the planning deviation isolates
how predictable the correspondence between scheduled and realised
durations is, regardless of the underlying procedure length.

Comparing both coefficients helps separate structural variability from
planning-related variability. The values in Table \ref{tab:cv_group}
show a clear scaling pattern in observed durations: relative variability
decreases as procedures become longer. Surgeries planned for less than
30 minutes have a CV of roughly 0.61, which falls steadily to 0.46 in
the 31--60 minute group, 0.39 in the 61--90 minute group, and stabilises
around 0.35 in the 91--180 minute range before rising slightly to 0.42
for procedures exceeding 180 minutes. Short procedures amplify small
absolute differences when expressed relative to the mean, while longer
procedures distribute variation across a larger time window.

The coefficients of variation for the planning deviation are
considerably higher in all groups, indicating a much broader
proportional spread in the alignment between planned and actual
durations. Values range from approximately 1.25 in the shortest category
to 1.71 in the 31--60 minute group, 1.06 in the 61--90 minute group,
0.91 in the 91--180 minute group, and 1.86 for procedures lasting more
than 180 minutes. This pattern suggests that longer and more complex
procedures, while not necessarily more variable in execution, show
greater uncertainty in planning accuracy.

Taken together, these patterns show that relative variability in
observed duration is predominantly structural and reflects differences
in case mix, whereas variability in planning deviation captures
fluctuations in scheduling precision. Distinguishing between these two
dimensions provides the analytical foundation for the next sections,
which examine how these dynamics unfold within each duration group and
across operating rooms.

\begin{table}

\caption{\label{tab:cv group}Summary of mean duration, SD, and CV by planned duration group (Observed vs. Planning deviation) \label{tab:cv_group}}
\centering
\begin{tabular}[t]{lrrrrrrr}
\toprule
Group & n & Mean duration & SD duration & CV duration & Mean planning & SD planning & CV planning\\
\midrule
<30 & 35 & 52.31429 & 31.96024 & 0.6109276 & 31.17143 & 27.77481 & 0.8910344\\
31–60 & 1673 & 60.84937 & 26.00058 & 0.4272941 & 16.68978 & 20.46855 & 1.2264126\\
61–90 & 3435 & 77.77671 & 30.58228 & 0.3932062 & 20.70830 & 21.67759 & 1.0468069\\
91–180 & 4018 & 112.94052 & 39.07618 & 0.3459890 & 24.89771 & 23.79474 & 0.9556997\\
> 180 & 121 & 168.28926 & 99.92326 & 0.5937590 & 356.78512 & 418.04701 & 1.1717053\\
\bottomrule
\end{tabular}
\end{table}

The monthly evolution of the coefficient of variation for observed
durations shows a clear and persistent separation between the planned
duration groups. The \textless30 minute category displays by far the
highest relative variability, with coefficients that fluctuate widely
from month to month and regularly exceed 0.7. This pronounced volatility
reflects the fact that even small absolute fluctuations translate into
large proportional changes for very short procedures.

In contrast, the 31--60, 61--90, and 91--180 minute groups follow stable
trajectories clustered between roughly 0.25 and 0.45 throughout the
observation period. Their lines move in parallel and show only modest
short-term variation, indicating that the proportional spread of
observed durations remains highly consistent within these mid-range
procedure lengths.

The \textgreater180 minute group occupies an intermediate position: its
coefficients are generally higher than those of the mid-duration groups
but lower and less erratic than the \textless30 group. Occasional peaks
appear, yet the overall pattern remains recognisable and bounded.

Taken together, the figure demonstrates that observed duration
variability scales predictably with procedure length. The short-duration
group exhibits structurally higher dispersion, whereas medium and
long-duration categories maintain a stable and well-defined variability
range across all months. Figure \ref{fig:cv_obs_group} illustrates these
patterns.

\begin{figure}[H]

{\centering \includegraphics[width=0.7\linewidth,]{In-Depth_Analysis_Cathlab_files/figure-latex/cv obs group-1} 

}

\caption{Monthly evolution of coefficient of variation (Observed durations) \label{fig:cv_obs_group}}\label{fig:cv obs group}
\end{figure}

The coefficient of variation of planning deviation displays
substantially more month-to-month fluctuation than the coefficient of
variation of observed durations, and the separation between duration
groups is less clear. The \textless30 minute group shows the most
volatile behaviour, with sharp spikes and dips that occasionally push
the coefficient well above 1.5 or below 0.5. Procedures lasting more
than 180 minutes also exhibit elevated variability, though their
fluctuations are less extreme than those of the shortest category. The
mid-range groups (31--60 and 61--90 minutes) follow a more irregular
pattern than might be expected, with values oscillating around 1.0 and
showing periodic peaks above 1.2. In contrast, the 91--180 minute group
remains the most stable, with coefficients clustering tightly around 0.9
to 1.1 across nearly all months. These patterns indicate that planning
accuracy is most consistent for medium-length procedures, while very
short and very long cases remain the most difficult to predict. Figure
\ref{fig:cv_planning_group} illustrates these monthly dynamics across
duration groups.

\begin{figure}[H]

{\centering \includegraphics[width=0.7\linewidth,]{In-Depth_Analysis_Cathlab_files/figure-latex/cv planning group-1} 

}

\caption{Monthly evolution of coefficient of variation (Planning deviation) \label{fig:cv_planning_group}}\label{fig:cv planning group}
\end{figure}

Viewed together, the two figures show that the relative variability of
observed surgery durations remains structurally stable across months and
across duration categories, with each group maintaining a consistent and
clearly separated pattern. By contrast, the variability in planning
deviation fluctuates much more strongly over time, particularly for the
shortest procedures and, to a lesser extent, the longest ones. These
contrasting patterns suggest that the execution of procedures is steady
and predictable, while planning accuracy is more sensitive to changes in
case mix, shifts in the share of complex cases, or month-to-month
variation in procedural demand. The subsequent sections examine these
dynamics in greater depth by analysing each duration group separately
and comparing how the coefficients of variation distribute across
operating rooms.

\vspace{1em}

\noindent\textbf{Surgeries planned under 30 minutes} \vspace{0.5em}

Short procedures, typically representing high volume and standardized
interventions, were intended to be analysed using coefficients of
variation at the operating room level. However, for this duration
category, no results are reported in Table \ref{tab:cv_room_lt30}. The
applied inclusion criterion, which requires a minimum number of
observations per operating room, was not met. As a result, the table
contains no entries.

\begin{longtable}[]{@{}
  >{\centering\arraybackslash}p{(\linewidth - 8\tabcolsep) * \real{0.1233}}
  >{\centering\arraybackslash}p{(\linewidth - 8\tabcolsep) * \real{0.1096}}
  >{\centering\arraybackslash}p{(\linewidth - 8\tabcolsep) * \real{0.2740}}
  >{\centering\arraybackslash}p{(\linewidth - 8\tabcolsep) * \real{0.2466}}
  >{\centering\arraybackslash}p{(\linewidth - 8\tabcolsep) * \real{0.2466}}@{}}
\caption{Coefficient of variation per operating room for surgeries
planned \textless30 minutes (only rooms with \textgreater10 surgeries
shown) \label{tab:cv_room_lt30} (continued below)}\tabularnewline
\toprule\noalign{}
\endfirsthead
\endhead
\bottomrule\noalign{}
\endlastfoot
\textbf{OR} & \textbf{n} & \textbf{Mean duration} & \textbf{SD duration}
& \textbf{CV duration} \\
\end{longtable}

\begin{longtable}[]{@{}
  >{\centering\arraybackslash}p{(\linewidth - 4\tabcolsep) * \real{0.2083}}
  >{\centering\arraybackslash}p{(\linewidth - 4\tabcolsep) * \real{0.1806}}
  >{\centering\arraybackslash}p{(\linewidth - 4\tabcolsep) * \real{0.1806}}@{}}
\toprule\noalign{}
\endhead
\bottomrule\noalign{}
\endlastfoot
\textbf{Mean dev} & \textbf{SD dev} & \textbf{CV dev} \\
\end{longtable}

\vspace{1em}

\noindent\textbf{Surgeries planned 31–60 minutes} \vspace{0.5em}

Procedures with a planned duration of 31--60 minutes show moderate
relative variability across operating rooms. For this category, the
coefficients of variation of observed durations range from 0.35 to 0.68,
as reported in Table \ref{tab:cv_room_31_60}. The lowest values appear
in KO01 (0.3544), KO02 (0.3583), and KO05 (0.3836), indicating
relatively homogeneous case durations in these rooms. Slightly higher
values are seen in KO07 (0.3902) and KO06 (0.4062), while KO03 (0.6755)
shows the greatest dispersion, reflecting a more heterogeneous mix of
procedures or the presence of atypical cases.

For the coefficient of variation of planning deviation, the values are
substantially higher across all rooms, ranging from 0.98 to 1.15. The
lowest dispersion appears in KO07 (0.9828) and KO02 (0.9918), followed
by KO03 (0.9993). Higher coefficients are observed in KO01 (1.1137),
KO06 (1.1524), and KO05 (1.1554), indicating greater variability in how
well planned times align with observed durations.

Taken together, these results show that surgeries planned for 31--60
minutes exhibit relatively consistent execution times but more variable
planning accuracy across operating rooms. The observed coefficients
largely reflect proportional scaling typical of medium-length procedures
rather than substantial operational differences between rooms.

\begin{longtable}[]{@{}
  >{\centering\arraybackslash}p{(\linewidth - 14\tabcolsep) * \real{0.0814}}
  >{\centering\arraybackslash}p{(\linewidth - 14\tabcolsep) * \real{0.0698}}
  >{\centering\arraybackslash}p{(\linewidth - 14\tabcolsep) * \real{0.1860}}
  >{\centering\arraybackslash}p{(\linewidth - 14\tabcolsep) * \real{0.1628}}
  >{\centering\arraybackslash}p{(\linewidth - 14\tabcolsep) * \real{0.1628}}
  >{\centering\arraybackslash}p{(\linewidth - 14\tabcolsep) * \real{0.1279}}
  >{\centering\arraybackslash}p{(\linewidth - 14\tabcolsep) * \real{0.1047}}
  >{\centering\arraybackslash}p{(\linewidth - 14\tabcolsep) * \real{0.1047}}@{}}
\caption{Coefficient of variation per operating room for surgeries
planned 31--60 minutes (only rooms with \textgreater10 surgeries shown)
\label{tab:cv_room_31_60}}\tabularnewline
\toprule\noalign{}
\begin{minipage}[b]{\linewidth}\centering
OR
\end{minipage} & \begin{minipage}[b]{\linewidth}\centering
n
\end{minipage} & \begin{minipage}[b]{\linewidth}\centering
Mean duration
\end{minipage} & \begin{minipage}[b]{\linewidth}\centering
SD duration
\end{minipage} & \begin{minipage}[b]{\linewidth}\centering
CV duration
\end{minipage} & \begin{minipage}[b]{\linewidth}\centering
Mean dev
\end{minipage} & \begin{minipage}[b]{\linewidth}\centering
SD dev
\end{minipage} & \begin{minipage}[b]{\linewidth}\centering
CV dev
\end{minipage} \\
\midrule\noalign{}
\endfirsthead
\toprule\noalign{}
\begin{minipage}[b]{\linewidth}\centering
OR
\end{minipage} & \begin{minipage}[b]{\linewidth}\centering
n
\end{minipage} & \begin{minipage}[b]{\linewidth}\centering
Mean duration
\end{minipage} & \begin{minipage}[b]{\linewidth}\centering
SD duration
\end{minipage} & \begin{minipage}[b]{\linewidth}\centering
CV duration
\end{minipage} & \begin{minipage}[b]{\linewidth}\centering
Mean dev
\end{minipage} & \begin{minipage}[b]{\linewidth}\centering
SD dev
\end{minipage} & \begin{minipage}[b]{\linewidth}\centering
CV dev
\end{minipage} \\
\midrule\noalign{}
\endhead
\bottomrule\noalign{}
\endlastfoot
KO02 & 535 & 55.77 & 21.38 & 0.3833 & 13.62 & 15 & 1.101 \\
KO05 & 450 & 58.96 & 22.62 & 0.3836 & 15.03 & 17.35 & 1.154 \\
KO01 & 363 & 56.1 & 20 & 0.3564 & 12.5 & 14.78 & 1.182 \\
KO06 & 242 & 78.5 & 32.52 & 0.4143 & 28.45 & 27.94 & 0.982 \\
KO07 & 70 & 70.04 & 37.28 & 0.5323 & 25.37 & 33.17 & 1.307 \\
KO03 & 11 & 96.18 & 64.98 & 0.6755 & 57.18 & 57.14 & 0.9993 \\
\end{longtable}

\vspace{1em}

\noindent\textbf{Surgeries planned 61–90 minutes} \vspace{0.5em}

Procedures planned between 61 and 90 minutes form a stable segment of
the surgical workflow, combining moderately complex interventions with
routine elective activity. Within this duration group, the coefficient
of variation of observed durations remains low and consistent across
operating rooms. As shown in Table \ref{tab:cv_room_61_90}, CV values
range from 0.3191 to 0.4616, indicating that execution times scale
proportionally and remain predictable once procedures begin.

The lowest variability appears in KO02 (0.3191) and KO01 (0.3708),
followed by KO05 (0.3781). Slightly higher dispersion is observed in
KO07 (0.4221), while the largest spread occurs in KO06 (0.4616), likely
reflecting a broader mix of procedure types or occasional atypical
cases. Despite these differences, the overall variation remains modest,
suggesting that mid-length operations follow consistent workflow and
anaesthesia routines.

For the coefficient of variation of planning deviation, the values are
higher but still concentrated within a moderate band. CV\_dev ranges
from 0.7459 to 1.24, with the lowest levels observed in KO02 (0.7459)
and KO01 (0.9448). Intermediate variability appears in KO05 (1.006) and
KO07 (1.075), while KO06 (1.24) displays the largest relative spread,
reflecting a wider range of deviations between planned and observed
timing.

Overall, surgeries in the 61--90 minute category show stable execution
times and moderate variation in planning accuracy across operating
rooms. Differences are primarily driven by case mix rather than
structural workflow variation, making this duration segment one of the
more predictable components of the surgical schedule.

\begin{longtable}[]{@{}
  >{\centering\arraybackslash}p{(\linewidth - 14\tabcolsep) * \real{0.0805}}
  >{\centering\arraybackslash}p{(\linewidth - 14\tabcolsep) * \real{0.0805}}
  >{\centering\arraybackslash}p{(\linewidth - 14\tabcolsep) * \real{0.1839}}
  >{\centering\arraybackslash}p{(\linewidth - 14\tabcolsep) * \real{0.1609}}
  >{\centering\arraybackslash}p{(\linewidth - 14\tabcolsep) * \real{0.1609}}
  >{\centering\arraybackslash}p{(\linewidth - 14\tabcolsep) * \real{0.1264}}
  >{\centering\arraybackslash}p{(\linewidth - 14\tabcolsep) * \real{0.1034}}
  >{\centering\arraybackslash}p{(\linewidth - 14\tabcolsep) * \real{0.1034}}@{}}
\caption{Coefficient of variation per operating room for surgeries
planned 61--90 minutes (only rooms with \textgreater10 surgeries shown)
\label{tab:cv_room_61_90}}\tabularnewline
\toprule\noalign{}
\begin{minipage}[b]{\linewidth}\centering
OR
\end{minipage} & \begin{minipage}[b]{\linewidth}\centering
n
\end{minipage} & \begin{minipage}[b]{\linewidth}\centering
Mean duration
\end{minipage} & \begin{minipage}[b]{\linewidth}\centering
SD duration
\end{minipage} & \begin{minipage}[b]{\linewidth}\centering
CV duration
\end{minipage} & \begin{minipage}[b]{\linewidth}\centering
Mean dev
\end{minipage} & \begin{minipage}[b]{\linewidth}\centering
SD dev
\end{minipage} & \begin{minipage}[b]{\linewidth}\centering
CV dev
\end{minipage} \\
\midrule\noalign{}
\endfirsthead
\toprule\noalign{}
\begin{minipage}[b]{\linewidth}\centering
OR
\end{minipage} & \begin{minipage}[b]{\linewidth}\centering
n
\end{minipage} & \begin{minipage}[b]{\linewidth}\centering
Mean duration
\end{minipage} & \begin{minipage}[b]{\linewidth}\centering
SD duration
\end{minipage} & \begin{minipage}[b]{\linewidth}\centering
CV duration
\end{minipage} & \begin{minipage}[b]{\linewidth}\centering
Mean dev
\end{minipage} & \begin{minipage}[b]{\linewidth}\centering
SD dev
\end{minipage} & \begin{minipage}[b]{\linewidth}\centering
CV dev
\end{minipage} \\
\midrule\noalign{}
\endhead
\bottomrule\noalign{}
\endlastfoot
KO05 & 1613 & 79.12 & 29.91 & 0.3781 & 20.59 & 20.72 & 1.006 \\
KO02 & 793 & 71.58 & 22.84 & 0.3191 & 18 & 13.43 & 0.7459 \\
KO06 & 454 & 89.07 & 41.12 & 0.4616 & 27.43 & 34.01 & 1.24 \\
KO01 & 357 & 72.55 & 26.9 & 0.3708 & 17.75 & 16.77 & 0.9448 \\
KO07 & 206 & 75.21 & 31.75 & 0.4221 & 21.36 & 22.97 & 1.075 \\
\end{longtable}

\vspace{1em}

\noindent\textbf{Surgeries planned 91–180 minutes} \vspace{0.5em}

Surgeries with a planned duration between 91 and 180 minutes include a
broad mix of major inpatient operations and complex elective cases.
Within this category, the coefficient of variation of observed durations
is moderate and shows limited dispersion across operating rooms. As
reported in Table \ref{tab:cv_room_91_180}, CV values range from 0.293
to 0.417, indicating proportionally consistent execution times despite
the wider absolute duration range.

The lowest variability is observed in KO02 (0.293) and KO01 (0.2932),
followed by KO07 (0.3283). Slightly higher values appear in KO06
(0.382), while the highest variability within this group is seen in KO05
(0.4175). These differences remain modest and largely reflect case mix
variation rather than structural operational differences.

The coefficient of variation of planning deviation follows a similarly
contained pattern. CV\_dev values range from 0.7761 to 1.003, with the
lowest variability observed in KO06 (0.7761) and KO01 (0.8813).
Intermediate dispersion appears in KO05 (0.8728) and KO02 (0.9077),
while the highest value is found in KO07 (1.003). These levels indicate
that planning accuracy for longer procedures is relatively stable, with
only moderate variation across operating rooms.

Together, the two indicators show that timing performance in the 91--180
minute category is consistent across rooms. Relative variability in
observed duration remains low, and variation in planning deviation is
moderate but tightly clustered. Compared with shorter procedures,
proportional deviations are smaller relative to total duration, making
these mid- to long-duration surgeries one of the most predictable
segments of the surgical schedule.

\begin{longtable}[]{@{}
  >{\centering\arraybackslash}p{(\linewidth - 14\tabcolsep) * \real{0.0805}}
  >{\centering\arraybackslash}p{(\linewidth - 14\tabcolsep) * \real{0.0805}}
  >{\centering\arraybackslash}p{(\linewidth - 14\tabcolsep) * \real{0.1839}}
  >{\centering\arraybackslash}p{(\linewidth - 14\tabcolsep) * \real{0.1609}}
  >{\centering\arraybackslash}p{(\linewidth - 14\tabcolsep) * \real{0.1609}}
  >{\centering\arraybackslash}p{(\linewidth - 14\tabcolsep) * \real{0.1264}}
  >{\centering\arraybackslash}p{(\linewidth - 14\tabcolsep) * \real{0.1034}}
  >{\centering\arraybackslash}p{(\linewidth - 14\tabcolsep) * \real{0.1034}}@{}}
\caption{Coefficient of variation per operating room for surgeries
planned 91--180 minutes (only rooms with \textgreater10 surgeries shown)
\label{tab:cv_room_91_180}}\tabularnewline
\toprule\noalign{}
\begin{minipage}[b]{\linewidth}\centering
OR
\end{minipage} & \begin{minipage}[b]{\linewidth}\centering
n
\end{minipage} & \begin{minipage}[b]{\linewidth}\centering
Mean duration
\end{minipage} & \begin{minipage}[b]{\linewidth}\centering
SD duration
\end{minipage} & \begin{minipage}[b]{\linewidth}\centering
CV duration
\end{minipage} & \begin{minipage}[b]{\linewidth}\centering
Mean dev
\end{minipage} & \begin{minipage}[b]{\linewidth}\centering
SD dev
\end{minipage} & \begin{minipage}[b]{\linewidth}\centering
CV dev
\end{minipage} \\
\midrule\noalign{}
\endfirsthead
\toprule\noalign{}
\begin{minipage}[b]{\linewidth}\centering
OR
\end{minipage} & \begin{minipage}[b]{\linewidth}\centering
n
\end{minipage} & \begin{minipage}[b]{\linewidth}\centering
Mean duration
\end{minipage} & \begin{minipage}[b]{\linewidth}\centering
SD duration
\end{minipage} & \begin{minipage}[b]{\linewidth}\centering
CV duration
\end{minipage} & \begin{minipage}[b]{\linewidth}\centering
Mean dev
\end{minipage} & \begin{minipage}[b]{\linewidth}\centering
SD dev
\end{minipage} & \begin{minipage}[b]{\linewidth}\centering
CV dev
\end{minipage} \\
\midrule\noalign{}
\endhead
\bottomrule\noalign{}
\endlastfoot
KO02 & 1510 & 108 & 31.66 & 0.293 & 21.12 & 19.17 & 0.9077 \\
KO01 & 1104 & 105.7 & 31 & 0.2932 & 21.43 & 18.88 & 0.8813 \\
KO07 & 630 & 142.4 & 46.75 & 0.3283 & 33.5 & 33.59 & 1.003 \\
KO05 & 534 & 104.1 & 43.45 & 0.4175 & 27.45 & 23.96 & 0.8728 \\
KO06 & 234 & 119.5 & 45.64 & 0.382 & 36.91 & 28.65 & 0.7761 \\
\end{longtable}

\vspace{1em}

\noindent\textbf{Surgeries planned over 180 minutes} \vspace{0.5em}

Procedures with a planned duration exceeding 180 minutes represent the
most complex and resource-intensive part of the surgical program. Within
this group, the coefficient of variation of observed durations shows a
wider spread across operating rooms than in the shorter duration
categories. As reported in Table \ref{tab:cv_room_over180}, CV values
range from 0.3288 to 0.795, indicating that despite the long procedure
times, relative variability remains moderate.

The lowest CV value is found in KO01 (0.3288), while KO05 (0.5518)
displays a noticeably higher proportional spread. The highest observed
CV\_duration occurs in KO06 (0.795), which reflects the substantial
heterogeneity and case-mix complexity associated with the limited number
of very long procedures performed in that room.

The coefficient of variation of planning deviation shows even greater
dispersion. Values range from 0.4157 up to 1.707, indicating that
proportional variation in planning accuracy is considerably larger than
the variation in observed duration itself. The lowest CV\_dev appears in
KO06 (0.4157), despite this room having the highest variability in
observed duration, while the highest value is recorded in KO01 (1.707).
KO05 shows a similarly elevated level of variability (1.508), consistent
with the broad spread in planning deviations for long and complex
surgeries.

Overall, operations lasting more than three hours exhibit moderate
proportional variability in actual duration but substantially higher
proportional variability in planning deviation. These patterns reflect
the intrinsic uncertainty of long and technically demanding procedures.
The room-to-room differences observed in this category are strongly
influenced by small sample sizes and case mix effects rather than
systematic operational inconsistencies.

\begin{longtable}[]{@{}
  >{\centering\arraybackslash}p{(\linewidth - 14\tabcolsep) * \real{0.0824}}
  >{\centering\arraybackslash}p{(\linewidth - 14\tabcolsep) * \real{0.0588}}
  >{\centering\arraybackslash}p{(\linewidth - 14\tabcolsep) * \real{0.1882}}
  >{\centering\arraybackslash}p{(\linewidth - 14\tabcolsep) * \real{0.1647}}
  >{\centering\arraybackslash}p{(\linewidth - 14\tabcolsep) * \real{0.1647}}
  >{\centering\arraybackslash}p{(\linewidth - 14\tabcolsep) * \real{0.1294}}
  >{\centering\arraybackslash}p{(\linewidth - 14\tabcolsep) * \real{0.1059}}
  >{\centering\arraybackslash}p{(\linewidth - 14\tabcolsep) * \real{0.1059}}@{}}
\caption{Coefficient of variation per operating room for surgeries
planned \textgreater{} 180 minutes (only rooms with \textgreater10
surgeries shown) \label{tab:cv_room_over180}}\tabularnewline
\toprule\noalign{}
\begin{minipage}[b]{\linewidth}\centering
OR
\end{minipage} & \begin{minipage}[b]{\linewidth}\centering
n
\end{minipage} & \begin{minipage}[b]{\linewidth}\centering
Mean duration
\end{minipage} & \begin{minipage}[b]{\linewidth}\centering
SD duration
\end{minipage} & \begin{minipage}[b]{\linewidth}\centering
CV duration
\end{minipage} & \begin{minipage}[b]{\linewidth}\centering
Mean dev
\end{minipage} & \begin{minipage}[b]{\linewidth}\centering
SD dev
\end{minipage} & \begin{minipage}[b]{\linewidth}\centering
CV dev
\end{minipage} \\
\midrule\noalign{}
\endfirsthead
\toprule\noalign{}
\begin{minipage}[b]{\linewidth}\centering
OR
\end{minipage} & \begin{minipage}[b]{\linewidth}\centering
n
\end{minipage} & \begin{minipage}[b]{\linewidth}\centering
Mean duration
\end{minipage} & \begin{minipage}[b]{\linewidth}\centering
SD duration
\end{minipage} & \begin{minipage}[b]{\linewidth}\centering
CV duration
\end{minipage} & \begin{minipage}[b]{\linewidth}\centering
Mean dev
\end{minipage} & \begin{minipage}[b]{\linewidth}\centering
SD dev
\end{minipage} & \begin{minipage}[b]{\linewidth}\centering
CV dev
\end{minipage} \\
\midrule\noalign{}
\endhead
\bottomrule\noalign{}
\endlastfoot
KO05 & 65 & 196.7 & 108.5 & 0.5518 & 226.7 & 341.7 & 1.508 \\
KO06 & 22 & 114.2 & 90.78 & 0.795 & 827.5 & 344 & 0.4157 \\
KO01 & 19 & 157.7 & 51.87 & 0.3288 & 172.7 & 294.8 & 1.707 \\
\end{longtable}

\subsubsection{Relationship between variability and mean operative
duration}\label{relationship-between-variability-and-mean-operative-duration}

Examining how variability scales with operative duration provides
insight into whether dispersion in surgical times increases
proportionally with the length of a procedure. By analysing this
relationship at the level of procedure types, it becomes possible to
distinguish structural scaling effects from differences in planning
performance. The following figures present three complementary
perspectives: the standard deviation of operative duration, the
coefficient of variation of observed duration, and the coefficient of
variation of planning deviation. Together, they describe how both
absolute and relative variability evolve as surgeries become longer or
more complex.

The first figure demonstrates a clear positive linear relationship
between mean operative time and its standard deviation, as shown in
Figure \ref{fig:sd_mean_relation}. Procedure types with longer average
durations also exhibit larger absolute variation, and this increase
occurs in a nearly proportional fashion. Shorter procedures with mean
durations around fifty to one hundred minutes show standard deviations
concentrated between twenty and fifty minutes, whereas procedures
exceeding two hundred minutes can display variability well above eighty
minutes. This scaling behaviour reflects the natural expansion of
absolute spread as operative time increases and does not imply a loss of
relative predictability. Instead, the pattern indicates that
proportional variability remains stable across procedure types even as
absolute deviations grow with case length.

\begin{figure}[H]

{\centering \includegraphics[width=0.7\linewidth,]{In-Depth_Analysis_Cathlab_files/figure-latex/sd mean relation-1} 

}

\caption{Standard deviation in operative time relative to mean duration \label{fig:sd_mean_relation}}\label{fig:sd mean relation}
\end{figure}

When variability is expressed relative to mean duration, a distinct
inverse relationship emerges, as shown in Figure
\ref{fig:cv_mean_relation}. The coefficient of variation declines
steadily as procedures become longer, indicating that proportional
dispersion shrinks with increasing operative time. Shorter procedures
frequently exhibit CV values between 0.5 and 0.6, reflecting their high
relative sensitivity to small absolute fluctuations, whereas most
operations exceeding one hundred minutes fall below 0.35. This pattern
demonstrates a fundamental scaling effect: although long procedures vary
more in absolute terms, their duration grows at a faster rate than their
spread. As a result, longer cases appear proportionally more stable,
while short procedures amplify even modest deviations. This inverse
gradient highlights the structural nature of temporal variability across
procedure types rather than differences in execution quality.

\begin{figure}[H]

{\centering \includegraphics[width=0.7\linewidth,]{In-Depth_Analysis_Cathlab_files/figure-latex/cv mean relation-1} 

}

\caption{Coefficient of variation (operative duration) relative to mean duration \label{fig:cv_mean_relation}}\label{fig:cv mean relation}
\end{figure}

A similar inverse relationship appears when examining planning deviation
relative to mean operative time, as shown in Figure
\ref{fig:cv_planning_relation}. Short procedures exhibit the widest
proportional spread, with coefficients of variation frequently ranging
between 1.5 and 3.0, and occasionally higher. These elevated values
arise because even modest absolute deviations translate into large
relative differences when planned durations are short. As procedures
become longer, this proportional inflation diminishes: CV values
gradually descend toward 1.0 and, for many longer procedures, fall below
that level. The downward slope in the figure therefore reflects a
scaling effect rather than a systematic change in planning performance.
Longer procedures have larger absolute deviations, but these are small
relative to their total planned duration, producing lower coefficients
of variation.

\begin{figure}[H]

{\centering \includegraphics[width=0.7\linewidth,]{In-Depth_Analysis_Cathlab_files/figure-latex/cv planning relation-1} 

}

\caption{Coefficient of variation (planning deviation) relative to mean duration \label{fig:cv_planning_relation}}\label{fig:cv planning relation}
\end{figure}

Taken together, the three figures show that while absolute variability
in surgical time increases with case length, relative variability tends
to decrease. This pattern indicates that proportional uncertainty is
greatest for short procedures, where small timing differences represent
a large share of the total duration, and least for long surgeries, where
absolute deviations are diluted across extended operating periods.

\subsection{How large are start-time delays, and when do they
happen?}\label{how-large-are-start-time-delays-and-when-do-they-happen}

Start-time deviations show a wide and asymmetric spread, with late
starts occurring more frequently and with larger magnitudes than early
starts. As shown in Table \ref{tab:start_diff}, 58.3 percent of
surgeries begin later than planned, with a mean delay of 69.6 minutes
and a median delay of 28 minutes. Early starts account for the remaining
41.7 percent of cases, averaging 56 minutes ahead of schedule with a
median of 31 minutes. Variability is substantial in both directions: the
standard deviation for late starts reaches 121 minutes, compared with 78
minutes for early starts, and the extreme upper tail extends to delays
of more than 1,400 minutes.

These patterns indicate that late starts represent the dominant source
of start-time deviation and pose the greatest operational risk, as
substantial delays early in the day can propagate through subsequent
cases and reduce overall room efficiency.

\begin{longtable}[]{@{}
  >{\centering\arraybackslash}p{(\linewidth - 20\tabcolsep) * \real{0.1573}}
  >{\centering\arraybackslash}p{(\linewidth - 20\tabcolsep) * \real{0.0787}}
  >{\centering\arraybackslash}p{(\linewidth - 20\tabcolsep) * \real{0.1011}}
  >{\centering\arraybackslash}p{(\linewidth - 20\tabcolsep) * \real{0.0899}}
  >{\centering\arraybackslash}p{(\linewidth - 20\tabcolsep) * \real{0.0674}}
  >{\centering\arraybackslash}p{(\linewidth - 20\tabcolsep) * \real{0.0674}}
  >{\centering\arraybackslash}p{(\linewidth - 20\tabcolsep) * \real{0.0787}}
  >{\centering\arraybackslash}p{(\linewidth - 20\tabcolsep) * \real{0.0899}}
  >{\centering\arraybackslash}p{(\linewidth - 20\tabcolsep) * \real{0.0674}}
  >{\centering\arraybackslash}p{(\linewidth - 20\tabcolsep) * \real{0.0899}}
  >{\centering\arraybackslash}p{(\linewidth - 20\tabcolsep) * \real{0.1124}}@{}}
\caption{Difference between actual and planned start time (absolute
minutes), separated by direction \label{tab:start_diff}}\tabularnewline
\toprule\noalign{}
\begin{minipage}[b]{\linewidth}\centering
direction
\end{minipage} & \begin{minipage}[b]{\linewidth}\centering
Mean
\end{minipage} & \begin{minipage}[b]{\linewidth}\centering
Median
\end{minipage} & \begin{minipage}[b]{\linewidth}\centering
SD
\end{minipage} & \begin{minipage}[b]{\linewidth}\centering
IQR
\end{minipage} & \begin{minipage}[b]{\linewidth}\centering
Min
\end{minipage} & \begin{minipage}[b]{\linewidth}\centering
Max
\end{minipage} & \begin{minipage}[b]{\linewidth}\centering
P95
\end{minipage} & \begin{minipage}[b]{\linewidth}\centering
P99
\end{minipage} & \begin{minipage}[b]{\linewidth}\centering
P99.9
\end{minipage} & \begin{minipage}[b]{\linewidth}\centering
Percent
\end{minipage} \\
\midrule\noalign{}
\endfirsthead
\toprule\noalign{}
\begin{minipage}[b]{\linewidth}\centering
direction
\end{minipage} & \begin{minipage}[b]{\linewidth}\centering
Mean
\end{minipage} & \begin{minipage}[b]{\linewidth}\centering
Median
\end{minipage} & \begin{minipage}[b]{\linewidth}\centering
SD
\end{minipage} & \begin{minipage}[b]{\linewidth}\centering
IQR
\end{minipage} & \begin{minipage}[b]{\linewidth}\centering
Min
\end{minipage} & \begin{minipage}[b]{\linewidth}\centering
Max
\end{minipage} & \begin{minipage}[b]{\linewidth}\centering
P95
\end{minipage} & \begin{minipage}[b]{\linewidth}\centering
P99
\end{minipage} & \begin{minipage}[b]{\linewidth}\centering
P99.9
\end{minipage} & \begin{minipage}[b]{\linewidth}\centering
Percent
\end{minipage} \\
\midrule\noalign{}
\endhead
\bottomrule\noalign{}
\endlastfoot
Late start & 69.6 & 28 & 121.3 & 73 & 1 & 1437 & 248 & 658 & 1181 &
58.3\% \\
Early start & 56 & 31 & 77.9 & 60 & 0 & 827 & 188.6 & 393 & 700.5 &
41.7\% \\
\end{longtable}

The distribution of start-time differences in Figure
\ref{fig:start_diff_hist} shows a strong concentration of cases near
zero, indicating that many surgeries begin close to their planned time.
The histogram, however, is clearly asymmetric: the bulk of activity is
shifted toward positive values, and the peak occurs slightly after zero,
reflecting the predominance of late starts over early ones.

Early starts do occur, but they form a thinner and shorter left-hand
tail, suggesting that substantial ahead-of-schedule starts are
relatively uncommon. In contrast, the right-hand tail is noticeably
longer and denser, revealing that while most delays remain modest, a
non-negligible subset of cases experiences large postponements. These
extended delays are infrequent but operationally significant, as they
can propagate through the day's schedule and reduce overall throughput.

Overall, the histogram illustrates a characteristic pattern of surgical
start-time behavior: many small deviations, a consistent bias toward
lateness, and occasional extreme delays that shape the right-skewed
distribution.

\begin{figure}[H]

{\centering \includegraphics[width=0.7\linewidth,]{In-Depth_Analysis_Cathlab_files/figure-latex/start diff hist-1} 

}

\caption{Start time deviations (grouped per 10 minutes) \label{fig:start_diff_hist}}\label{fig:start diff hist}
\end{figure}

Average start-time delays vary sharply by weekday, as shown in Figure
\ref{fig:start_delay_weekday}. During regular working days, start delays
remain low and tightly clustered, ranging from roughly 8 to 18 minutes.
Monday (11.5 minutes), Tuesday (14.1 minutes), Wednesday (8.3 minutes),
Thursday (18.5 minutes), and Friday (9.7 minutes) all follow a similar
pattern, indicating consistent and well-controlled morning starts
throughout the standard workweek.

A completely different picture emerges on weekends. Mean delays rise
dramatically to 355.8 minutes on Saturday and 343.6 minutes on Sunday,
an increase of several hours compared with weekday performance. These
extreme values reflect the nature of weekend activity, which is
typically dominated by urgent or unscheduled procedures rather than
structured elective programs. Such cases often occur later in the day,
outside the usual planned start windows, which inflates the recorded
deviation relative to scheduled time.

The contrast between weekdays and weekends is therefore clear: weekday
activity shows predictable, well-aligned start times, whereas weekend
operations display very long apparent delays that are better interpreted
as the result of irregular and demand-driven scheduling rather than
punctuality issues.

\begin{figure}[H]

{\centering \includegraphics[width=0.7\linewidth,]{In-Depth_Analysis_Cathlab_files/figure-latex/start delay weekday-1} 

}

\caption{Mean start delay by weekday \label{fig:start_delay_weekday}}\label{fig:start delay weekday}
\end{figure}

Late starts are more common than early starts on every weekday, as shown
in Table \ref{tab:start_diff_weekday}. Across Monday to Friday, between
56 and 59 percent of surgeries begin later than scheduled. Mean delays
range from 57.6 to 69.7 minutes, while the median delay stays
consistently near 24 to 30 minutes, indicating a stable pattern of
moderate lateness across the workweek. Early starts occur in about 39 to
41 percent of cases on weekdays, with mean early leads between 49 and 63
minutes and median leads close to −29 to −35 minutes. These small
differences across Monday through Friday suggest a largely uniform level
of punctuality during the standard operating schedule.

Weekend activity departs sharply from this pattern. On Saturday and
Sunday, over 86 percent of cases start late, and early starts drop to 8
to 14 percent. Mean delays rise dramatically to 434.5 minutes on both
days, with median delays of 354 to 384 minutes, reflecting the fact that
weekend procedures often occur far outside the planned daytime schedule.
Likewise, mean early starts show very large negative values (around −229
to −233 minutes), again illustrating the atypical timing of weekend
operations. These large deviations do not reflect punctuality in the
traditional sense but the fundamentally different scheduling dynamics of
urgent or ad hoc procedures performed outside standard program hours.

\begin{longtable}[]{@{}
  >{\centering\arraybackslash}p{(\linewidth - 14\tabcolsep) * \real{0.0926}}
  >{\centering\arraybackslash}p{(\linewidth - 14\tabcolsep) * \real{0.0648}}
  >{\centering\arraybackslash}p{(\linewidth - 14\tabcolsep) * \real{0.1667}}
  >{\centering\arraybackslash}p{(\linewidth - 14\tabcolsep) * \real{0.1204}}
  >{\centering\arraybackslash}p{(\linewidth - 14\tabcolsep) * \real{0.1389}}
  >{\centering\arraybackslash}p{(\linewidth - 14\tabcolsep) * \real{0.1759}}
  >{\centering\arraybackslash}p{(\linewidth - 14\tabcolsep) * \real{0.1111}}
  >{\centering\arraybackslash}p{(\linewidth - 14\tabcolsep) * \real{0.1296}}@{}}
\caption{Start-time deviation patterns per weekday (minutes).
\label{tab:start_diff_weekday}}\tabularnewline
\toprule\noalign{}
\begin{minipage}[b]{\linewidth}\centering
Weekday
\end{minipage} & \begin{minipage}[b]{\linewidth}\centering
n
\end{minipage} & \begin{minipage}[b]{\linewidth}\centering
Late starts (\%)
\end{minipage} & \begin{minipage}[b]{\linewidth}\centering
Mean delay
\end{minipage} & \begin{minipage}[b]{\linewidth}\centering
Median delay
\end{minipage} & \begin{minipage}[b]{\linewidth}\centering
Early starts (\%)
\end{minipage} & \begin{minipage}[b]{\linewidth}\centering
Mean lead
\end{minipage} & \begin{minipage}[b]{\linewidth}\centering
Median lead
\end{minipage} \\
\midrule\noalign{}
\endfirsthead
\toprule\noalign{}
\begin{minipage}[b]{\linewidth}\centering
Weekday
\end{minipage} & \begin{minipage}[b]{\linewidth}\centering
n
\end{minipage} & \begin{minipage}[b]{\linewidth}\centering
Late starts (\%)
\end{minipage} & \begin{minipage}[b]{\linewidth}\centering
Mean delay
\end{minipage} & \begin{minipage}[b]{\linewidth}\centering
Median delay
\end{minipage} & \begin{minipage}[b]{\linewidth}\centering
Early starts (\%)
\end{minipage} & \begin{minipage}[b]{\linewidth}\centering
Mean lead
\end{minipage} & \begin{minipage}[b]{\linewidth}\centering
Median lead
\end{minipage} \\
\midrule\noalign{}
\endhead
\bottomrule\noalign{}
\endlastfoot
Ma & 1398 & 56.2 & 64.3 & 29.5 & 41.4 & -59.5 & -35 \\
Di & 1812 & 58.6 & 57.6 & 24 & 39.6 & -49.6 & -29 \\
Wo & 2308 & 57.2 & 58 & 29 & 41.3 & -60.2 & -31 \\
Do & 2034 & 58.5 & 69.7 & 24 & 39.5 & -56.4 & -34 \\
Vr & 1598 & 58.9 & 59.2 & 28.5 & 39.4 & -63.9 & -39 \\
Za & 59 & 86.4 & 434.5 & 384 & 8.5 & -233.4 & -347 \\
Zo & 73 & 86.3 & 434.6 & 354 & 13.7 & -229.6 & -249.5 \\
\end{longtable}

Start-time reliability differs substantially across operating rooms, as
shown in Table \ref{tab:start_diff_or}. The proportion of late starts
ranges from 30 percent to 78 percent, indicating meaningful but not
extreme variation across the OR suite.

The highest shares of late starts appear in KO06 (78.4 percent) and KO07
(76.8 percent). These rooms also differ strongly in the magnitude of the
associated delays. KO06 stands out with a mean delay of 205 minutes and
a median of 127 minutes, suggesting that a small number of very late
cases heavily influence the average. By contrast, KO07, despite a
similarly high percentage of late starts, shows a much smaller mean
delay of 52.6 minutes and a median of 22 minutes, indicating more
moderate deviations.

Other large rooms, such as KO02 (57.3 percent late starts) and KO05
(50.5 percent late starts), show more balanced patterns. Mean delays in
these rooms lie around 44 to 50 minutes, with median delays near 20
minutes, suggesting relatively stable morning and intra-day starts with
occasional larger disruptions.

At the lower end, KO01 (51.8 percent late) and especially KO04 (30
percent late) demonstrate comparatively strong punctuality. KO04 shows
the smallest delays overall, with a mean delay of only 9 minutes and a
median of 3 minutes, reflecting highly predictable start-time
performance.

Early starts also vary across rooms. In KO05 and KO01, early starts
occur nearly as often as late ones (around 46--47 percent), with mean
leads between 30 and 52 minutes. In KO06, the mean early start of −135.9
minutes again highlights the presence of rare but extreme timing
irregularities.

Taken together, the results indicate that most operating rooms
experience relatively stable start-time performance, with moderate
delays and a balanced mix of early and late starts. A small number of
rooms---particularly KO06---drive much of the observed variability due
to infrequent but very large deviations. These differences suggest that
case mix, scheduling patterns, or workflow constraints may differ
significantly between rooms, shaping their punctuality profiles.

\begin{longtable}[]{@{}
  >{\centering\arraybackslash}p{(\linewidth - 14\tabcolsep) * \real{0.1009}}
  >{\centering\arraybackslash}p{(\linewidth - 14\tabcolsep) * \real{0.0642}}
  >{\centering\arraybackslash}p{(\linewidth - 14\tabcolsep) * \real{0.1651}}
  >{\centering\arraybackslash}p{(\linewidth - 14\tabcolsep) * \real{0.1193}}
  >{\centering\arraybackslash}p{(\linewidth - 14\tabcolsep) * \real{0.1376}}
  >{\centering\arraybackslash}p{(\linewidth - 14\tabcolsep) * \real{0.1743}}
  >{\centering\arraybackslash}p{(\linewidth - 14\tabcolsep) * \real{0.1101}}
  >{\centering\arraybackslash}p{(\linewidth - 14\tabcolsep) * \real{0.1284}}@{}}
\caption{Start-time deviation patterns per operating room (minutes).
\label{tab:start_diff_or}}\tabularnewline
\toprule\noalign{}
\begin{minipage}[b]{\linewidth}\centering
ActualOR
\end{minipage} & \begin{minipage}[b]{\linewidth}\centering
n
\end{minipage} & \begin{minipage}[b]{\linewidth}\centering
Late starts (\%)
\end{minipage} & \begin{minipage}[b]{\linewidth}\centering
Mean delay
\end{minipage} & \begin{minipage}[b]{\linewidth}\centering
Median delay
\end{minipage} & \begin{minipage}[b]{\linewidth}\centering
Early starts (\%)
\end{minipage} & \begin{minipage}[b]{\linewidth}\centering
Mean lead
\end{minipage} & \begin{minipage}[b]{\linewidth}\centering
Median lead
\end{minipage} \\
\midrule\noalign{}
\endfirsthead
\toprule\noalign{}
\begin{minipage}[b]{\linewidth}\centering
ActualOR
\end{minipage} & \begin{minipage}[b]{\linewidth}\centering
n
\end{minipage} & \begin{minipage}[b]{\linewidth}\centering
Late starts (\%)
\end{minipage} & \begin{minipage}[b]{\linewidth}\centering
Mean delay
\end{minipage} & \begin{minipage}[b]{\linewidth}\centering
Median delay
\end{minipage} & \begin{minipage}[b]{\linewidth}\centering
Early starts (\%)
\end{minipage} & \begin{minipage}[b]{\linewidth}\centering
Mean lead
\end{minipage} & \begin{minipage}[b]{\linewidth}\centering
Median lead
\end{minipage} \\
\midrule\noalign{}
\endhead
\bottomrule\noalign{}
\endlastfoot
KO06 & 961 & 78.4 & 205.2 & 127 & 20.7 & -135.9 & -76 \\
KO07 & 914 & 76.8 & 52.6 & 22 & 21.8 & -53.9 & -32 \\
KO03 & 34 & 58.8 & 203 & 59.5 & 35.3 & -62.9 & -50.5 \\
KO02 & 2852 & 57.3 & 44.1 & 23 & 40.6 & -47.5 & -30 \\
KO01 & 1846 & 51.8 & 43.4 & 19 & 46.6 & -52 & -31 \\
KO05 & 2665 & 50.5 & 50.2 & 21 & 47.4 & -61.5 & -35 \\
KO04 & 10 & 30 & 9 & 3 & 70 & -67.9 & -68 \\
\end{longtable}

\subsection{Which procedures deviate most from
plan?}\label{which-procedures-deviate-most-from-plan}

These procedures represent complex or less frequently performed
interventions, where case-to-case variability is high and precise
planning is therefore more difficult. In several of these procedures,
the share of longer-than-planned cases exceeds 65 percent, confirming a
consistent tendency toward underestimation.

More common procedures with moderately high deviations, such as RX
interventionele neuro (44.2 minutes) or RX interventionele perifeer
(43.5 minutes), show a more balanced split between longer and shorter
cases. Their deviations likely reflect broad case-mix variability within
the same procedural label rather than structural misplanning.

At the lower end of the top twenty list---but still with deviations
exceeding 22 minutes---procedures such as Percutane aortaklep
vervanging, BRONCHOSCOPIE+AFNAME, and E.R.C.P. show deviations driven by
heterogeneity in disease severity, anatomical complexity, or required
adjunct steps.

Across almost all listed procedures, both positive and negative
deviations occur, but the mean + deviation values tend to be larger in
magnitude than the mean -- deviation, indicating that
longer-than-planned cases contribute more to the overall deviation than
shorter ones. This pattern aligns with the operational reality that
unexpected complexity more often prolongs procedures than shortens them.

Overall, the results illustrate that procedure type is a strong
determinant of planning accuracy. Routine or standardized interventions
show predictable timing, while complex, technical, or less frequently
performed procedures exhibit substantial uncertainty, necessitating
generous scheduling buffers and careful resource allocation.

\begin{longtable}[]{@{}
  >{\centering\arraybackslash}p{(\linewidth - 12\tabcolsep) * \real{0.2640}}
  >{\centering\arraybackslash}p{(\linewidth - 12\tabcolsep) * \real{0.0480}}
  >{\centering\arraybackslash}p{(\linewidth - 12\tabcolsep) * \real{0.1680}}
  >{\centering\arraybackslash}p{(\linewidth - 12\tabcolsep) * \real{0.1040}}
  >{\centering\arraybackslash}p{(\linewidth - 12\tabcolsep) * \real{0.1520}}
  >{\centering\arraybackslash}p{(\linewidth - 12\tabcolsep) * \real{0.1120}}
  >{\centering\arraybackslash}p{(\linewidth - 12\tabcolsep) * \real{0.1520}}@{}}
\caption{Top 20 procedures by average absolute duration deviation
(Cathlab). \label{tab:duration_dev_proc}}\tabularnewline
\toprule\noalign{}
\begin{minipage}[b]{\linewidth}\centering
ProcedureName
\end{minipage} & \begin{minipage}[b]{\linewidth}\centering
n
\end{minipage} & \begin{minipage}[b]{\linewidth}\centering
Mean abs dev (min)
\end{minipage} & \begin{minipage}[b]{\linewidth}\centering
Longer (\%)
\end{minipage} & \begin{minipage}[b]{\linewidth}\centering
Mean + dev (min)
\end{minipage} & \begin{minipage}[b]{\linewidth}\centering
Shorter (\%)
\end{minipage} & \begin{minipage}[b]{\linewidth}\centering
Mean - dev (min)
\end{minipage} \\
\midrule\noalign{}
\endfirsthead
\toprule\noalign{}
\begin{minipage}[b]{\linewidth}\centering
ProcedureName
\end{minipage} & \begin{minipage}[b]{\linewidth}\centering
n
\end{minipage} & \begin{minipage}[b]{\linewidth}\centering
Mean abs dev (min)
\end{minipage} & \begin{minipage}[b]{\linewidth}\centering
Longer (\%)
\end{minipage} & \begin{minipage}[b]{\linewidth}\centering
Mean + dev (min)
\end{minipage} & \begin{minipage}[b]{\linewidth}\centering
Shorter (\%)
\end{minipage} & \begin{minipage}[b]{\linewidth}\centering
Mean - dev (min)
\end{minipage} \\
\midrule\noalign{}
\endhead
\bottomrule\noalign{}
\endlastfoot
Recanalisatie intracranieel stroke & 207 & 85.7 & 49.8 & 44.7 & 50.2 &
-126.2 \\
embolisatie visceraal & 99 & 61.2 & 69.7 & 35.4 & 30.3 & -120.5 \\
LEAD EXTRACTIE & 35 & 58.9 & 68.6 & 70.7 & 31.4 & -33 \\
ABLATIE STRUCTURELE VENTRICULAIRE TACHYCARDIE & 100 & 56.4 & 41 & 49.7 &
59 & -61.1 \\
RX interventionele neuro & 449 & 44.2 & 51.4 & 45.8 & 47.2 & -43.7 \\
RX interventionele perifeer & 90 & 43.5 & 55.6 & 28.1 & 44.4 & -62.8 \\
RX angio 4-vaten & 38 & 38.8 & 78.9 & 45.4 & 18.4 & -16.3 \\
ABLATIE ATRIALE TACHYCARDIE & 211 & 33 & 48.8 & 35.7 & 50.2 & -31 \\
ENDOPROTHESE & 119 & 31.2 & 52.1 & 37.5 & 44.5 & -26.3 \\
ABLATIE IDIOPATISCHE VENTRICULAIRE TACHYCARDIE & 213 & 30.5 & 44.6 & 35
& 54 & -27.5 \\
DILATATIE FEMORALIS & 818 & 29.6 & 43.6 & 25.3 & 55.1 & -33.6 \\
ABLATIE WPW & 96 & 28.9 & 57.3 & 36.2 & 41.7 & -19.8 \\
DILATATIE ILIACA KISSING & 142 & 26.7 & 40.8 & 28.8 & 57 & -26.3 \\
DILATATIE POPLITEA & 43 & 26.7 & 41.9 & 35.8 & 58.1 & -20.2 \\
ABLATIE PERSISTERENDE VKF & 155 & 26.4 & 36.8 & 21.9 & 62.6 & -29.4 \\
DILATATIE TIBIALIS & 273 & 24.1 & 47.3 & 30.7 & 50.9 & -18.9 \\
RX Vertebroplastie & 170 & 24.1 & 92.4 & 25.7 & 5.9 & -6.9 \\
E.R.C.P. & 60 & 23 & 18.3 & 22.5 & 80 & -23.6 \\
BRONCHOSCOPIE+AFNAME & 188 & 22.9 & 19.7 & 20.5 & 77.7 & -24.3 \\
Percutane aortaklep vervanging & 237 & 22.4 & 40.1 & 25 & 59.1 &
-20.9 \\
\end{longtable}

The scatter plot in Figure \ref{fig:volume_deviation} illustrates how
procedure volume relates to planning accuracy. A clear pattern emerges
at the lower end of the volume scale: rare procedures tend to show the
highest mean absolute deviations, with several low-volume interventions
exceeding 50 to 80 minutes of average deviation. This confirms that
infrequent or highly specialized procedures are more difficult to
estimate reliably due to limited historical data and greater
case-to-case variability.

As procedure volume increases, deviations become more compact, and
high-volume procedures cluster at lower mean deviations, generally below
25 minutes. This indicates that repetition and accumulated experience
improve the accuracy of duration estimates.

However, the relationship is not strictly linear. Among mid- and
high-volume procedures, the scatter widens again, and the downward trend
becomes less pronounced. Several procedures with substantial case counts
still exhibit deviations of 30 to 40 minutes, suggesting that factors
beyond volume---such as procedural heterogeneity or patient
complexity---also play a meaningful role.

Taken together, the figure shows that although volume is a strong
predictor of planning accuracy at the extremes, it only partially
explains variation in the middle range. High familiarity improves
predictability, but even frequently performed procedures can deviate
substantially when clinical complexity varies.

\begin{figure}[H]

{\centering \includegraphics[width=0.7\linewidth,]{In-Depth_Analysis_Cathlab_files/figure-latex/volume_deviation-1} 

}

\caption{Relationship between procedure volume and mean deviation \label{fig:volume_deviation}}\label{fig:volume_deviation}
\end{figure}

The distribution of duration deviations across surgeons shows
substantial heterogeneity in planning accuracy, as summarized in Table
\ref{tab:surgeon_deviation}. Among surgeons with at least fifty recorded
procedures, the mean absolute deviation ranges from roughly 20 to 78
minutes, indicating a wide spectrum of prediction performance. A small
subset of surgeons exhibit markedly higher deviations---above 50
minutes---while most fall between 23 and 40 minutes, suggesting that for
the majority, planning accuracy remains reasonably stable.

Relative deviations display a similarly broad pattern, ranging from
18.3\% to 43.2\%, with several surgeons exceeding 30\%, reflecting
procedures or surgeon-specific styles that introduce greater
variability. Importantly, the surgeon with the highest absolute
deviation (77.7 minutes) does not have the lowest case volume, and
surgeons with very large procedure counts (e.g., 931, 1169, 1025 cases)
show deviations comparable to or better than colleagues with far fewer
cases. This confirms that case volume alone does not explain the
differences in planning accuracy.

Overall, the results highlight that duration deviations are not
uniformly distributed across surgeons. Even within a shared operational
environment---and often within similar procedural
categories---individual practice patterns, case mix, and surgical
complexity contribute meaningfully to the observed variability.

\begin{longtable}[]{@{}lccc@{}}
\caption{Top 15 surgeons by average absolute and relative duration
deviation (pseudonymized). \label{tab:surgeon_deviation}}\tabularnewline
\toprule\noalign{}
Surgeon & n & Mean absolute dev (min) & Mean relative dev (\%) \\
\midrule\noalign{}
\endfirsthead
\toprule\noalign{}
Surgeon & n & Mean absolute dev (min) & Mean relative dev (\%) \\
\midrule\noalign{}
\endhead
\bottomrule\noalign{}
\endlastfoot
Surgeon 1 & 107 & 77.7 & 31.5 \\
Surgeon 2 & 465 & 51.7 & 37.4 \\
Surgeon 3 & 64 & 43.5 & 36.3 \\
Surgeon 4 & 454 & 43.5 & 38.8 \\
Surgeon 5 & 74 & 40.0 & 43.2 \\
Surgeon 6 & 266 & 28.5 & 22.0 \\
Surgeon 7 & 522 & 26.8 & 28.6 \\
Surgeon 8 & 1431 & 26.7 & 26.7 \\
Surgeon 9 & 177 & 25.3 & 25.4 \\
Surgeon 10 & 271 & 25.3 & 28.6 \\
Surgeon 11 & 173 & 24.5 & 43.8 \\
Surgeon 12 & 87 & 23.6 & 31.3 \\
Surgeon 13 & 232 & 23.6 & 29.5 \\
Surgeon 14 & 1169 & 21.0 & 21.3 \\
Surgeon 15 & 1025 & 20.1 & 18.3 \\
\end{longtable}

\newpage

\section{Do surgeries stay within regular working
hours?}\label{do-surgeries-stay-within-regular-working-hours}

\subsection{How often do cases extend beyond their assigned time
window?}\label{how-often-do-cases-extend-beyond-their-assigned-time-window}

Around 5.8 percent of all surgeries at Cathlab extend beyond the
scheduled shift end of the operating team, as shown in Table
\ref{tab:overtime_summary}. Among cases that run late, the average
overtime is 41.8 minutes, while the median is 32 minutes, indicating
that most overruns remain modest in length. The 95th percentile of 108.3
minutes shows that only a small share of procedures continue
substantially beyond the shift, contributing disproportionately to the
total amount of after-hours work.

These results suggest that overtime is a recurring but contained
component of the surgical schedule. While the vast majority of
procedures finish within the planned working period, the steady presence
of late-running cases highlights the need for operational buffers
between shifts. The distribution of overtime durations further indicates
that although extended sessions do occur, they are relatively infrequent
and typically concentrated among a small number of longer or more
unpredictable procedures.

If you want, I can now help you write the subsection that breaks this
down by weekday, operating room, urgency type, or procedure group.

\begin{longtable}[]{@{}
  >{\centering\arraybackslash}p{(\linewidth - 10\tabcolsep) * \real{0.0833}}
  >{\centering\arraybackslash}p{(\linewidth - 10\tabcolsep) * \real{0.2708}}
  >{\centering\arraybackslash}p{(\linewidth - 10\tabcolsep) * \real{0.1354}}
  >{\centering\arraybackslash}p{(\linewidth - 10\tabcolsep) * \real{0.1667}}
  >{\centering\arraybackslash}p{(\linewidth - 10\tabcolsep) * \real{0.1875}}
  >{\centering\arraybackslash}p{(\linewidth - 10\tabcolsep) * \real{0.1562}}@{}}
\caption{Overall share and duration of after-hours surgeries (Cathlab).
\label{tab:overtime_summary}}\tabularnewline
\toprule\noalign{}
\begin{minipage}[b]{\linewidth}\centering
Total
\end{minipage} & \begin{minipage}[b]{\linewidth}\centering
Surgeries with overtime
\end{minipage} & \begin{minipage}[b]{\linewidth}\centering
Percentage
\end{minipage} & \begin{minipage}[b]{\linewidth}\centering
Mean Overtime
\end{minipage} & \begin{minipage}[b]{\linewidth}\centering
Median Overtime
\end{minipage} & \begin{minipage}[b]{\linewidth}\centering
P95 Overtime
\end{minipage} \\
\midrule\noalign{}
\endfirsthead
\toprule\noalign{}
\begin{minipage}[b]{\linewidth}\centering
Total
\end{minipage} & \begin{minipage}[b]{\linewidth}\centering
Surgeries with overtime
\end{minipage} & \begin{minipage}[b]{\linewidth}\centering
Percentage
\end{minipage} & \begin{minipage}[b]{\linewidth}\centering
Mean Overtime
\end{minipage} & \begin{minipage}[b]{\linewidth}\centering
Median Overtime
\end{minipage} & \begin{minipage}[b]{\linewidth}\centering
P95 Overtime
\end{minipage} \\
\midrule\noalign{}
\endhead
\bottomrule\noalign{}
\endlastfoot
9282 & 539 & 5.8 & 41.8 & 32 & 108.3 \\
\end{longtable}

The distribution of overtime duration, shown in Figure
\ref{fig:overtime_dist}, is highly right-skewed. Most after-hours cases
conclude within the first hour beyond the scheduled shift, with the
largest concentration of surgeries falling between 0 and 40 minutes of
overtime. Beyond this range, frequencies decline steadily. Only a small
number of cases exceed 100 minutes, and instances above 200 or 300
minutes appear as isolated outliers.

This pattern indicates that while overtime occurs regularly, it is
usually modest in length. The long right tail highlights that a small
subset of prolonged procedures contributes disproportionately to total
overtime minutes. These uncommon but extended cases can place additional
pressure on staffing and late-day turnover, yet they do not reflect the
typical after-hours workload. Overall, overtime at Cathlab primarily
consists of short extensions to planned schedules rather than routine,
extensive overruns.

\begin{figure}[H]

{\centering \includegraphics[width=0.7\linewidth,]{In-Depth_Analysis_Cathlab_files/figure-latex/overtime_dist-1} 

}

\caption{Distribution of overtime duration for after-hours surgeries \label{fig:overtime_dist}}\label{fig:overtime_dist}
\end{figure}

he share of surgeries that continue beyond regular shift hours remains
low and relatively stable during the workweek, as shown in Figure
\ref{fig:afterhours_weekday}. From Monday to Friday, after-hours
activity fluctuates between 3.5 and 7.1 percent, indicating that the
vast majority of weekday procedures finish within planned operating
windows. The small variation between days suggests that weekday
scheduling is well contained, with only a modest portion of the workload
extending past shift boundaries.

A pronounced shift appears during the weekend. On Saturday, 20.3 percent
of surgeries run beyond scheduled hours. The highest share of the week
and on Sunday the proportion remains elevated at 15.1 percent. Although
weekend surgical volume is substantially lower, these higher percentages
signal that cases performed on Saturdays and Sundays are more likely to
exceed their allocated time. This pattern reflects the dominant presence
of urgent or semi-urgent procedures on weekends, which are inherently
less predictable in both timing and duration.

Taken together, the contrast between weekdays and weekends shows that
after-hours work at Cathlab is structurally linked to emergency-driven
activity rather than routine scheduling. Weekday overtime remains
limited and controlled, while weekend operations more often spill over
into the next shift due to the unpredictable nature of the case mix.

\begin{figure}[H]

{\centering \includegraphics[width=0.7\linewidth,]{In-Depth_Analysis_Cathlab_files/figure-latex/afterhours_weekday-1} 

}

\caption{Share of after-hours surgeries by weekday \label{fig:afterhours_weekday}}\label{fig:afterhours_weekday}
\end{figure}

The proportion of surgeries extending beyond scheduled shift hours
varies modestly from year to year, as shown in Figure
\ref{fig:afterhours_trend}. In 2022, 6.0 percent of all surgeries at
Cathlab continued past their planned shift end. This share rose to 6.9
percent in 2023 before dropping noticeably to 4.7 percent in 2024. The
partial data for 2025 show an increase to 5.8 percent, returning closer
to earlier levels but still below the peak observed in 2023.

These values represent the percentage of all procedures in each year
that required overtime. Rather than a consistent upward or downward
trajectory, the pattern shows alternating increases and decreases,
suggesting that year-to-year variation is likely driven by shifts in
case mix, throughput pressure, or episodic operational factors rather
than structural change.

Overall, overtime remains a stable feature of the operating schedule,
with annual shares clustering between roughly 5 and 7 percent. The
temporary dip in 2024 indicates that overtime can fluctuate meaningfully
from one year to the next, but no sustained trend toward either
improvement or deterioration is visible over the observed period.

\begin{figure}[H]

{\centering \includegraphics[width=0.7\linewidth,]{In-Depth_Analysis_Cathlab_files/figure-latex/afterhours_trend-1} 

}

\caption{Evolution of after-hours activity over time \label{fig:afterhours_trend}}\label{fig:afterhours_trend}
\end{figure}

The proportion of surgeries extending beyond scheduled shift hours
varies substantially across operating rooms, as shown in Table
\ref{tab:overtime_or}. KO03 exhibits the highest overtime share at 17.6
percent, more than triple that of several other rooms. Although KO03
handles a relatively small number of cases, its overtime episodes are
long, with a mean of 65 minutes and a P95 of 206.8 minutes, indicating
occasional substantial spillovers.

KO06 and KO04 follow with overtime rates of 12.8 and 10 percent,
respectively. KO06 in particular shows sizeable extensions, with a
median overtime of 40 minutes and a P95 of 118.2 minutes. These patterns
suggest that workflows in these rooms regularly encounter timing
pressures, likely due to case complexity or variability in procedure
duration.

The remaining rooms, KO07, KO02, KO01, and KO05, display more moderate
overtime frequencies between 4.5 and 6.7 percent. Their average overtime
durations fall between 31.6 and 55.6 minutes, and their medians between
24 and 46 minutes. These values indicate that while overtime does occur,
it is generally contained and less frequent than in the higher-pressure
rooms.

Overall, overtime activity is unevenly distributed across the operating
suite. A small number of rooms account for a disproportionate share of
after-hours work, driven by longer or more unpredictable procedures.
Most rooms, however, maintain relatively stable and predictable
schedules, with overtime occurring only occasionally and typically of
moderate duration.

\begin{longtable}[]{@{}
  >{\centering\arraybackslash}p{(\linewidth - 10\tabcolsep) * \real{0.1279}}
  >{\centering\arraybackslash}p{(\linewidth - 10\tabcolsep) * \real{0.0930}}
  >{\centering\arraybackslash}p{(\linewidth - 10\tabcolsep) * \real{0.2907}}
  >{\centering\arraybackslash}p{(\linewidth - 10\tabcolsep) * \real{0.1860}}
  >{\centering\arraybackslash}p{(\linewidth - 10\tabcolsep) * \real{0.2093}}
  >{\centering\arraybackslash}p{(\linewidth - 10\tabcolsep) * \real{0.0930}}@{}}
\caption{Operating rooms by after-hours activity (Cathlab).
\label{tab:overtime_or}}\tabularnewline
\toprule\noalign{}
\begin{minipage}[b]{\linewidth}\centering
ActualOR
\end{minipage} & \begin{minipage}[b]{\linewidth}\centering
Cases
\end{minipage} & \begin{minipage}[b]{\linewidth}\centering
Percentage of overtime surgeries
\end{minipage} & \begin{minipage}[b]{\linewidth}\centering
Mean Overtime
\end{minipage} & \begin{minipage}[b]{\linewidth}\centering
Median Overtime
\end{minipage} & \begin{minipage}[b]{\linewidth}\centering
P95
\end{minipage} \\
\midrule\noalign{}
\endfirsthead
\toprule\noalign{}
\begin{minipage}[b]{\linewidth}\centering
ActualOR
\end{minipage} & \begin{minipage}[b]{\linewidth}\centering
Cases
\end{minipage} & \begin{minipage}[b]{\linewidth}\centering
Percentage of overtime surgeries
\end{minipage} & \begin{minipage}[b]{\linewidth}\centering
Mean Overtime
\end{minipage} & \begin{minipage}[b]{\linewidth}\centering
Median Overtime
\end{minipage} & \begin{minipage}[b]{\linewidth}\centering
P95
\end{minipage} \\
\midrule\noalign{}
\endhead
\bottomrule\noalign{}
\endlastfoot
KO03 & 34 & 17.6 & 65 & 31 & 206.8 \\
KO06 & 961 & 12.8 & 50.2 & 40 & 118.2 \\
KO04 & 10 & 10 & 40 & 40 & 40 \\
KO07 & 914 & 6.7 & 55.6 & 46 & 180 \\
KO02 & 2852 & 4.9 & 31.6 & 28 & 72 \\
KO01 & 1846 & 4.7 & 32.9 & 24 & 89 \\
KO05 & 2665 & 4.5 & 43.3 & 33 & 116.2 \\
\end{longtable}

\subsection{When does overtime occur, and how heavy is
it?}\label{when-does-overtime-occur-and-how-heavy-is-it}

The distribution of overtime surgeries, shown in Figure
\ref{fig:overtime_timing}, is dominated by cases that finish shortly
after 16:30, the official end of the day shift. The highest bar appears
immediately after this threshold, indicating that many procedures begun
late in the afternoon run only modestly beyond scheduled hours. A second
noticeable cluster extends into the early evening, with rapidly
decreasing frequencies toward 18:00 and 19:00.

Beyond this early-evening window, the number of cases drops sharply.
Only a small fraction of surgeries end between 20:00 and 22:00, and
activity becomes even more sparse thereafter. Surgeries finishing after
midnight appear rarely, and only a handful extend into the early morning
hours, with isolated cases recorded as late as 08:30.

Overall, the timing pattern shows that overtime at Cathlab is driven
mainly by moderate extensions of the daytime program rather than by
persistent late-night operations. The long tail into the evening and
night reflects infrequent but occasionally lengthy procedures, while the
bulk of overtime surgeries conclude within one to two hours after the
regular shift ends.

\begin{figure}[H]

{\centering \includegraphics[width=0.7\linewidth,]{In-Depth_Analysis_Cathlab_files/figure-latex/overtime_timing-1} 

}

\caption{Timing of overtime surgeries \label{fig:overtime_timing}}\label{fig:overtime_timing}
\end{figure}

Average overtime durations vary moderately across weekdays, as shown in
Figure \ref{fig:overtime_weekday}. From Monday to Friday, mean overtime
ranges between roughly 33 and 43 minutes, while the corresponding
medians fall between about 25 and 32 minutes. The consistent gap between
mean and median confirms a right-skewed pattern, in which most overruns
are modest but a smaller number of longer cases pull the mean upward.

Weekend cases show noticeably larger extensions. On Saturday, the mean
overtime rises to about 60 minutes with a median of 38 minutes, and on
Sunday the mean reaches more than 70 minutes with a median of 41
minutes. These higher values reflect the different case mix on weekends,
where a greater share of time-unpredictable urgent procedures is
performed.

Across the full week, a clear pattern emerges: when overtime occurs, the
typical extension (as reflected by medians) lies around 25--40 minutes
on weekdays and slightly above 40 minutes on weekends, while the average
extension is consistently higher due to occasional long-running cases.
This combination of stable weekday patterns and more variable weekend
behaviour highlights the influence of case mix on overtime duration
rather than systematic differences in day-to-day performance.

\begin{figure}[H]

{\centering \includegraphics[width=0.7\linewidth,]{In-Depth_Analysis_Cathlab_files/figure-latex/overtime_weekday-1} 

}

\caption{Average and median overtime duration by weekday \label{fig:overtime_weekday}}\label{fig:overtime_weekday}
\end{figure}

Average and median overtime durations per surgery show a modest degree
of year-to-year fluctuation, as illustrated in Figure
\ref{fig:overtime_year}, but no structural upward or downward drift. The
mean overtime decreases from 45.9 minutes in 2022 to 42.6 minutes in
2023, reaches its lowest level in 2024 at 35.1 minutes, and rises again
to 43.3 minutes in 2025. The median follows a similar pattern, declining
from 34 minutes in 2022 to 29 minutes in 2024, before increasing to 41
minutes in 2025.

Across all years, medians remain well below the means, confirming a
right-skewed overtime distribution in which most overruns are modest
while a smaller number of prolonged cases elevates the average. The
typical overrun (median) ranges from roughly 29 to 41 minutes, while the
mean hovers around 40--46 minutes.

Overall, the pattern suggests stable overtime behaviour with normal
operational variation rather than systematic improvement or
deterioration. The persistent gap between mean and median highlights
that a limited subset of long-running procedures continues to account
for a disproportionate share of total overtime minutes, even as the bulk
of cases finish within a relatively narrow and predictable range of
additional time.

\begin{figure}[H]

{\centering \includegraphics[width=0.7\linewidth,]{In-Depth_Analysis_Cathlab_files/figure-latex/overtime_year-1} 

}

\caption{Average and median overtime per surgery by year \label{fig:overtime_year}}\label{fig:overtime_year}
\end{figure}

The distribution of overtime varies strongly by shift type, as
illustrated in Figure \ref{fig:overtime_shift}. The day shift
(08:00--16:30) shows by far the greatest spread, with a dense
concentration of short extensions near zero but a long right tail
reaching well above 150 minutes. This indicates that while most
day-shift surgeries overrun by only a modest amount, a small number of
prolonged cases contribute disproportionately to total overtime. The day
shift therefore represents the primary source of extended after-hours
activity simply because it handles the highest volume and the widest mix
of procedures.

In contrast, overtime during the evening shift (16:30--22:00) is more
narrowly concentrated. The majority of cases cluster around relatively
small extensions, and only a few observations depart substantially from
the shift boundary. The more compact shape of the distribution suggests
that evening activity, although occasionally prolonged, is generally
well contained.

The night shift (22:00--08:00) displays an even tighter distribution,
reflecting the limited number of surgeries performed during these hours.
Most procedures end close to the scheduled time, and only a handful show
large deviations. The low density in the violin plot corresponds
directly to the smaller case volume overnight.

Taken together, the three distributions show a clear pattern: overtime
is structurally most variable during daytime hours, when throughput is
highest and the likelihood of spillover rises with case complexity.
Evening and night shifts exhibit shorter, more predictable overruns,
with only sporadic long extensions. Across all shifts, skewed right
tails confirm that a small number of unusually long cases continue to
account for a substantial share of total overtime minutes.

\begin{figure}[H]

{\centering \includegraphics[width=0.7\linewidth,]{In-Depth_Analysis_Cathlab_files/figure-latex/overtime_shift-1} 

}

\caption{Relative overtime per shift type \label{fig:overtime_shift}}\label{fig:overtime_shift}
\end{figure}

\subsection{How does Cathlab compare with Genk, Lanaken, and
Maaseik?}\label{how-does-cathlab-compare-with-genk-lanaken-and-maaseik}

The comparison between campuses shows marked differences in the extent
and intensity of after-hours surgical activity, as summarised in Table
\ref{tab:afterhours_campus}. Genk stands out with the highest share of
overtime surgeries at 8.4 percent, confirming that most extended
operations occur at the main campus. Cathlab follows with a moderate
share of 5.8 percent, clearly above Maaseik (2.8 percent) and Lanaken
(0.5 percent). This suggests that Cathlab maintains a higher level of
late-day intervention than the smaller campuses, likely due to its
specific procedural focus and the nature of cardiovascular work that
often runs beyond regular hours.

The average overtime at Cathlab is 41.8 minutes, which is shorter than
Genk's 59 minutes but comparable to Maaseik's 42 minutes, and notably
higher than Lanaken's 15.7 minutes. Median overtime shows a similar
pattern, with 32 minutes at Cathlab versus 38 minutes at Genk, 29
minutes at Maaseik, and 10 minutes at Lanaken.

Overall, the results indicate that while Genk dominates in total volume
and absolute duration of after-hours work, Cathlab experiences a
proportionally high share of extended cases relative to its size. This
highlights Cathlab's intermediate operational profile, less intensive
than Genk but more time-sensitive and variable than the smaller
peripheral sites.

\begin{longtable}[]{@{}
  >{\centering\arraybackslash}p{(\linewidth - 12\tabcolsep) * \real{0.0820}}
  >{\centering\arraybackslash}p{(\linewidth - 12\tabcolsep) * \real{0.0656}}
  >{\centering\arraybackslash}p{(\linewidth - 12\tabcolsep) * \real{0.1803}}
  >{\centering\arraybackslash}p{(\linewidth - 12\tabcolsep) * \real{0.2459}}
  >{\centering\arraybackslash}p{(\linewidth - 12\tabcolsep) * \real{0.1557}}
  >{\centering\arraybackslash}p{(\linewidth - 12\tabcolsep) * \real{0.1475}}
  >{\centering\arraybackslash}p{(\linewidth - 12\tabcolsep) * \real{0.1230}}@{}}
\caption{Comparison of after-hours activity between campuses.
\label{tab:afterhours_campus}}\tabularnewline
\toprule\noalign{}
\begin{minipage}[b]{\linewidth}\centering
Campus
\end{minipage} & \begin{minipage}[b]{\linewidth}\centering
Total
\end{minipage} & \begin{minipage}[b]{\linewidth}\centering
AfterHour surgeries
\end{minipage} & \begin{minipage}[b]{\linewidth}\centering
Share of overtime surgeries
\end{minipage} & \begin{minipage}[b]{\linewidth}\centering
Average Overtime
\end{minipage} & \begin{minipage}[b]{\linewidth}\centering
Median Overtime
\end{minipage} & \begin{minipage}[b]{\linewidth}\centering
P95 Overtime
\end{minipage} \\
\midrule\noalign{}
\endfirsthead
\toprule\noalign{}
\begin{minipage}[b]{\linewidth}\centering
Campus
\end{minipage} & \begin{minipage}[b]{\linewidth}\centering
Total
\end{minipage} & \begin{minipage}[b]{\linewidth}\centering
AfterHour surgeries
\end{minipage} & \begin{minipage}[b]{\linewidth}\centering
Share of overtime surgeries
\end{minipage} & \begin{minipage}[b]{\linewidth}\centering
Average Overtime
\end{minipage} & \begin{minipage}[b]{\linewidth}\centering
Median Overtime
\end{minipage} & \begin{minipage}[b]{\linewidth}\centering
P95 Overtime
\end{minipage} \\
\midrule\noalign{}
\endhead
\bottomrule\noalign{}
\endlastfoot
Genk & 96044 & 8024 & 8.4 & 59 & 38 & 193.8 \\
Cathlab & 9282 & 539 & 5.8 & 41.8 & 32 & 108.3 \\
Maaseik & 53902 & 1494 & 2.8 & 42 & 29 & 129 \\
Lanaken & 69395 & 371 & 0.5 & 15.7 & 10 & 48.5 \\
\end{longtable}

\newpage

\section{Are surgeries performed in their assigned operating
rooms?}\label{are-surgeries-performed-in-their-assigned-operating-rooms}

\subsection{How frequently are surgeries relocated to a different
room?}\label{how-frequently-are-surgeries-relocated-to-a-different-room}

Room reallocations are relatively rare at Cathlab, as shown in Table
\ref{tab:room_swap}. Out of 9282 recorded surgeries, only 113 were
performed in a different operating room than originally planned,
corresponding to 1.2 percent of all cases. This very low share indicates
that room assignments are followed with a high degree of consistency,
and that unexpected relocations are the exception rather than the rule.

Although infrequent, these reallocations remain operationally meaningful
because they usually arise from reactive adjustments to delays,
equipment needs, or urgent case insertions. Even a small number of such
changes can generate additional coordination work for staff, influence
room preparation timing, and affect turnover efficiency. Monitoring
their occurrence therefore offers insight into the stability of the
scheduling process and helps determine whether deviations are sporadic
or linked to particular rooms, time periods, or operational conditions.

\begin{longtable}[]{@{}
  >{\centering\arraybackslash}p{(\linewidth - 4\tabcolsep) * \real{0.1111}}
  >{\centering\arraybackslash}p{(\linewidth - 4\tabcolsep) * \real{0.2917}}
  >{\centering\arraybackslash}p{(\linewidth - 4\tabcolsep) * \real{0.1944}}@{}}
\caption{Overall frequency of surgeries relocated to a different room
(Cathlab). \label{tab:room_swap}}\tabularnewline
\toprule\noalign{}
\begin{minipage}[b]{\linewidth}\centering
Total
\end{minipage} & \begin{minipage}[b]{\linewidth}\centering
Swapped (absolute)
\end{minipage} & \begin{minipage}[b]{\linewidth}\centering
Swapped (\%)
\end{minipage} \\
\midrule\noalign{}
\endfirsthead
\toprule\noalign{}
\begin{minipage}[b]{\linewidth}\centering
Total
\end{minipage} & \begin{minipage}[b]{\linewidth}\centering
Swapped (absolute)
\end{minipage} & \begin{minipage}[b]{\linewidth}\centering
Swapped (\%)
\end{minipage} \\
\midrule\noalign{}
\endhead
\bottomrule\noalign{}
\endlastfoot
9282 & 113 & 1.2 \\
\end{longtable}

The share of relocated surgeries remains very low across all weekdays,
as shown in Figure \ref{fig:room_swap_weekday}, with values generally
fluctuating between 0.3 and 2.5 percent. This consistency indicates that
room assignments are highly stable during regular working days, and that
unplanned reallocations are only rarely required. Monday and Wednesday
show the lowest relocation rates, while Tuesday and Thursday reach
slightly higher levels, though still well below three percent.

A noticeable increase appears during the weekend. On Saturday, 6.9
percent of surgeries are performed in a different room than planned, and
on Sunday the share is 1.4 percent. These higher weekend values likely
reflect the reduced number of available operating rooms and the less
predictable mix of urgent or ad hoc cases, which can require occasional
reactive adjustments to room allocation.

\begin{figure}[H]

{\centering \includegraphics[width=0.7\linewidth,]{In-Depth_Analysis_Cathlab_files/figure-latex/room_swap_weekday-1} 

}

\caption{Share of relocated surgeries by weekday \label{fig:room_swap_weekday}}\label{fig:room_swap_weekday}
\end{figure}

The monthly percentage of surgeries relocated to a different operating
room shows substantial short-term fluctuation, ranging from close to
zero to more than 3 percent, as seen in Figure
\ref{fig:room_swap_month}. Several months exhibit brief spikes above 2
percent, while others fall well below 1 percent. These swings indicate
that room reallocations occur irregularly and reflect short-lived
operational pressures rather than persistent patterns.

Despite this variability, no clear upward or downward trend emerges
across the observed period. The baseline remains low throughout, and the
high points appear as isolated events rather than evidence of structural
instability in the planning process. Overall, the monthly pattern
confirms that room swaps are rare, reactive adjustments that occur only
when specific scheduling or logistical constraints arise, while the
underlying room assignment process remains steady over time.

\begin{figure}[H]

{\centering \includegraphics[width=0.7\linewidth,]{In-Depth_Analysis_Cathlab_files/figure-latex/room_swap_month-1} 

}

\caption{Monthly percentage of surgeries with room swap \label{fig:room_swap_month}}\label{fig:room_swap_month}
\end{figure}

\subsection{Are specific rooms affected more than
others?}\label{are-specific-rooms-affected-more-than-others}

The analysis shows that only a small subset of operating rooms is
regularly involved in relocations, as summarized in Table
\ref{tab:room_swap_or}. Most planned rooms record only one or two
reallocations across the entire period, while a few display recurring
patterns. The largest number of moved cases originates from KO01, KO06,
and KO02, which together account for the majority of all swaps. These
rooms consistently redirect procedures to a limited set of destination
rooms, most notably KO02 and KO07, suggesting that these receiving rooms
serve as flexible buffers when minor scheduling conflicts arise.

The relocation patterns are highly structured. For example, KO01 sends
more than 90 percent of its relocations to KO02, while KO06 redirects
nearly all moved cases to KO07. KO02 shows a similar dominant pairing
with KO01. These concentrated flows indicate that swaps are not random
adjustments but follow stable operational relationships between rooms
that share similar equipment, staffing profiles, or procedural
portfolios.

Rooms with very low volumes of relocated cases, such as KO03, KO05,
GO04, GO11, and GO08, show no meaningful secondary patterns, and the
rare swaps originating from GOR0 and KO04 all move to a single
consistent destination. Overall, the data suggest that room
reallocations are infrequent and, when they do occur, follow predictable
and functionally sensible pathways within the operating complex.

\begin{longtable}[]{@{}
  >{\centering\arraybackslash}p{(\linewidth - 8\tabcolsep) * \real{0.1500}}
  >{\centering\arraybackslash}p{(\linewidth - 8\tabcolsep) * \real{0.1625}}
  >{\centering\arraybackslash}p{(\linewidth - 8\tabcolsep) * \real{0.2375}}
  >{\centering\arraybackslash}p{(\linewidth - 8\tabcolsep) * \real{0.2250}}
  >{\centering\arraybackslash}p{(\linewidth - 8\tabcolsep) * \real{0.2250}}@{}}
\caption{Top 10 planned ORs with the highest number of room swaps,
showing the five most frequent new OR destinations (Cathlab).
\label{tab:room_swap_or}}\tabularnewline
\toprule\noalign{}
\begin{minipage}[b]{\linewidth}\centering
PlannedOR
\end{minipage} & \begin{minipage}[b]{\linewidth}\centering
TotalMoved
\end{minipage} & \begin{minipage}[b]{\linewidth}\centering
Top1
\end{minipage} & \begin{minipage}[b]{\linewidth}\centering
Top2
\end{minipage} & \begin{minipage}[b]{\linewidth}\centering
Top3
\end{minipage} \\
\midrule\noalign{}
\endfirsthead
\toprule\noalign{}
\begin{minipage}[b]{\linewidth}\centering
PlannedOR
\end{minipage} & \begin{minipage}[b]{\linewidth}\centering
TotalMoved
\end{minipage} & \begin{minipage}[b]{\linewidth}\centering
Top1
\end{minipage} & \begin{minipage}[b]{\linewidth}\centering
Top2
\end{minipage} & \begin{minipage}[b]{\linewidth}\centering
Top3
\end{minipage} \\
\midrule\noalign{}
\endhead
\bottomrule\noalign{}
\endlastfoot
KO01 & 32 & KO02 (30, 93.8\%) & KO03 (2, 6.2\%) & NA \\
KO06 & 28 & KO07 (25, 89.3\%) & KO05 (3, 10.7\%) & NA \\
KO02 & 21 & KO01 (19, 90.5\%) & KO03 (1, 4.8\%) & KO05 (1, 4.8\%) \\
KO03 & 4 & KO02 (3, 75\%) & KO07 (1, 25\%) & NA \\
KO05 & 3 & KO02 (1, 33.3\%) & KO03 (1, 33.3\%) & KO07 (1, 33.3\%) \\
GO04 & 2 & KO04 (1, 50\%) & KO05 (1, 50\%) & NA \\
GO11 & 2 & KO06 (2, 100\%) & NA & NA \\
GO08 & 1 & KO05 (1, 100\%) & NA & NA \\
GOR0 & 1 & KO07 (1, 100\%) & NA & NA \\
KO04 & 1 & KO05 (1, 100\%) & NA & NA \\
\end{longtable}

A small number of operating rooms receive the majority of relocated
surgeries, as shown in Table \ref{tab:room_swap_dest}. KO02 is the most
common destination, receiving 34 relocated cases, followed by KO07 with
28 cases and KO01 with 20. Together, these three rooms account for a
clear majority of all relocations. KO06 also appears frequently with 19
received cases, while KO05, KO03, and KO04 receive only a small number
of transfers.

The pattern is highly concentrated. Most relocations are absorbed by
just four rooms: KO02, KO07, KO01, and KO06. However, the remaining ORs
receive transfers only occasionally. This suggests that these rooms have
operational characteristics that make them suitable fallback locations,
such as broader procedural compatibility, more flexible staffing
availability, or favorable scheduling profiles that allow last-minute
adjustments.

Although relocations are rare overall, the fact that they consistently
funnel into the same small set of rooms indicates a structured and
intentional use of capacity buffers within the operating complex.

\begin{longtable}[]{@{}
  >{\centering\arraybackslash}p{(\linewidth - 4\tabcolsep) * \real{0.1528}}
  >{\centering\arraybackslash}p{(\linewidth - 4\tabcolsep) * \real{0.2222}}
  >{\centering\arraybackslash}p{(\linewidth - 4\tabcolsep) * \real{0.1528}}@{}}
\caption{Top 15 actual operating rooms most frequently used as
destinations after relocation (Cathlab).
\label{tab:room_swap_dest}}\tabularnewline
\toprule\noalign{}
\begin{minipage}[b]{\linewidth}\centering
ActualOR
\end{minipage} & \begin{minipage}[b]{\linewidth}\centering
ReceivedCases
\end{minipage} & \begin{minipage}[b]{\linewidth}\centering
SharePct
\end{minipage} \\
\midrule\noalign{}
\endfirsthead
\toprule\noalign{}
\begin{minipage}[b]{\linewidth}\centering
ActualOR
\end{minipage} & \begin{minipage}[b]{\linewidth}\centering
ReceivedCases
\end{minipage} & \begin{minipage}[b]{\linewidth}\centering
SharePct
\end{minipage} \\
\midrule\noalign{}
\endhead
\bottomrule\noalign{}
\endlastfoot
KO02 & 34 & 0.37 \\
KO07 & 28 & 0.3 \\
KO01 & 20 & 0.22 \\
KO06 & 19 & 0.2 \\
KO05 & 7 & 0.08 \\
KO03 & 4 & 0.04 \\
KO04 & 1 & 0.01 \\
\end{longtable}

Procedures with the highest likelihood of being relocated include
several specialized cardiology and radiology interventions, as shown in
Table \ref{tab:room_swap_procedure}. The procedures at the top of the
list---such as RX angio 4-vaten, pacemaker implantation, lead
extraction, and various neuro- and peripheral radiological interventions
show relocation rates ranging from about 3 to 11 percent, well above the
overall relocation rate of roughly 1 percent.

What these procedures have in common is not low complexity, but high
substitutability of room infrastructure. Many of them can be performed
in multiple hybrid or radiology-equipped operating rooms, which makes
them more likely to be moved when scheduling constraints arise. Their
technical requirements are specialized but not restricted to a single
room, giving planners flexibility when balancing resource availability
across the day.

Lower in the ranking we also find several electrophysiology procedures
(for example ABLATE WPW, ABLATE BUNDEL HIS, and various atrial
tachycardia ablations), along with selected neurological and
gastrointestinal interventions. These cases typically require particular
equipment but can still be executed in more than one compatible
environment, explaining their elevated swap percentages.

By contrast, large surgical specialties with highly dedicated
infrastructure---orthopaedics, major abdominal surgery, obstetrics---are
virtually absent from the list, reflecting minimal relocation for cases
that depend on room-specific setups.

Overall, the table shows that room relocations at Cathlab are
concentrated among procedure types that are technically specialised yet
operationally flexible, allowing them to be reassigned between multiple
suitable rooms without compromising clinical or logistical requirements.

\begin{longtable}[]{@{}
  >{\centering\arraybackslash}p{(\linewidth - 4\tabcolsep) * \real{0.4167}}
  >{\centering\arraybackslash}p{(\linewidth - 4\tabcolsep) * \real{0.0833}}
  >{\centering\arraybackslash}p{(\linewidth - 4\tabcolsep) * \real{0.1389}}@{}}
\caption{Top 15 procedure types with the highest share of relocated
surgeries (Cathlab). \label{tab:room_swap_procedure}}\tabularnewline
\toprule\noalign{}
\begin{minipage}[b]{\linewidth}\centering
ProcedureName
\end{minipage} & \begin{minipage}[b]{\linewidth}\centering
n
\end{minipage} & \begin{minipage}[b]{\linewidth}\centering
SwapPct
\end{minipage} \\
\midrule\noalign{}
\endfirsthead
\toprule\noalign{}
\begin{minipage}[b]{\linewidth}\centering
ProcedureName
\end{minipage} & \begin{minipage}[b]{\linewidth}\centering
n
\end{minipage} & \begin{minipage}[b]{\linewidth}\centering
SwapPct
\end{minipage} \\
\midrule\noalign{}
\endhead
\bottomrule\noalign{}
\endlastfoot
RX angio 4-vaten & 38 & 10.81 \\
PACEMAKER PLAATSING & 81 & 8.64 \\
LEAD EXTRACTIE & 35 & 5.71 \\
RX interventionele neuro & 449 & 4.46 \\
RX interventionele perifeer & 90 & 4.44 \\
Recanalisatie intracranieel stroke & 207 & 3.86 \\
ABLATIE WPW & 96 & 3.12 \\
PTC + dilatatie & 99 & 3.03 \\
embolisatie visceraal & 99 & 3.03 \\
AICD IMPLANTATIE & 39 & 2.56 \\
ABLATIE BUNDEL HIS & 237 & 1.69 \\
E.R.C.P. & 60 & 1.67 \\
ABLATIE ATRIALE TACHYCARDIE & 211 & 1.42 \\
ABLATIE IDIOPATISCHE VENTRICULAIRE TACHYCARDIE & 213 & 1.41 \\
ABLATIE AV NODALE REENTRY TACHYCARDIE & 350 & 1.14 \\
\end{longtable}

\subsection{What is the operational impact of these
reallocations?}\label{what-is-the-operational-impact-of-these-reallocations}

Relocated surgeries differ from non-relocated ones in several dimensions
of start-time behaviour, but they show broadly similar patterns in
duration deviations and overtime, as summarized in Table
\ref{tab:room_swap_impact}.

For start times, relocated cases exhibit substantially larger delays
among late starts. The mean late-start deviation for relocated surgeries
is 125.3 minutes compared with 68.8 minutes for non-relocated cases. The
median values show a similar but less extreme pattern (42 versus 28
minutes). This indicates that when a relocated case begins later than
planned, the delay tends to be much more pronounced, consistent with the
added coordination required to prepare a different operating room than
originally scheduled.

A similar effect appears for early starts, where relocated surgeries
start further ahead of schedule on average (72.6 versus 58.4 minutes).
Medians again are closer (39 versus 33 minutes), suggesting that while
most relocated cases differ modestly from plan, a limited number
experience large shifts in timing that inflate the mean.

In contrast, deviations in operative duration show no meaningful
difference between relocated and non-relocated procedures.
Longer-than-planned cases diverge by roughly 23 to 27 minutes on average
in both groups, and shorter-than-planned cases show nearly identical
medians of 18 to 19 minutes. This indicates that once a surgery begins,
its execution time is not affected by whether it was moved to another
room.

Overtime is rare in both groups. The median is zero minutes for
relocated and non-relocated surgeries, and the mean overtime is only
slightly higher for relocated cases (4.8 versus 2.4 minutes). This
confirms that relocations do not typically cause surgeries to run into
the next shift.

Taken together, the results suggest that the operational impact of
relocation is concentrated at the start of the surgical process, where
setup changes and room availability constraints lead to larger
deviations from the planned start time. Once underway, however,
relocated surgeries behave similarly to non-relocated ones in duration
accuracy and overtime exposure.

\begin{longtable}[]{@{}
  >{\centering\arraybackslash}p{(\linewidth - 6\tabcolsep) * \real{0.2222}}
  >{\centering\arraybackslash}p{(\linewidth - 6\tabcolsep) * \real{0.2639}}
  >{\centering\arraybackslash}p{(\linewidth - 6\tabcolsep) * \real{0.1111}}
  >{\centering\arraybackslash}p{(\linewidth - 6\tabcolsep) * \real{0.1250}}@{}}
\caption{Comparison of start and duration deviations (early/late,
shorter/longer) and overtime between relocated and non-relocated
surgeries (Cathlab). \label{tab:room_swap_impact}}\tabularnewline
\toprule\noalign{}
\begin{minipage}[b]{\linewidth}\centering
Type
\end{minipage} & \begin{minipage}[b]{\linewidth}\centering
Metric
\end{minipage} & \begin{minipage}[b]{\linewidth}\centering
Mean
\end{minipage} & \begin{minipage}[b]{\linewidth}\centering
Median
\end{minipage} \\
\midrule\noalign{}
\endfirsthead
\toprule\noalign{}
\begin{minipage}[b]{\linewidth}\centering
Type
\end{minipage} & \begin{minipage}[b]{\linewidth}\centering
Metric
\end{minipage} & \begin{minipage}[b]{\linewidth}\centering
Mean
\end{minipage} & \begin{minipage}[b]{\linewidth}\centering
Median
\end{minipage} \\
\midrule\noalign{}
\endhead
\bottomrule\noalign{}
\endlastfoot
Relocated & Late start & 125.3 & 42 \\
Not relocated & Late start & 68.8 & 28 \\
& ---- & & \\
Relocated & Early start & 72.6 & 39 \\
Not relocated & Early start & 58.4 & 33 \\
& ---- & & \\
Relocated & Longer duration & 23.4 & 18 \\
Not relocated & Longer duration & 26.9 & 18 \\
& ---- & & \\
Relocated & Shorter duration & 83.8 & 19 \\
Not relocated & Shorter duration & 26 & 16 \\
& ---- & & \\
Relocated & Overtime & 4.8 & 0 \\
Not relocated & Overtime & 2.4 & 0 \\
& ---- & & \\
\end{longtable}

Overtime is somewhat more common among relocated surgeries than among
those performed in their originally assigned operating room, as shown in
Table \ref{tab:overtime_relocation}. Among non relocated cases, 5.7
percent extend into after hours, whereas the share rises to 10.6 percent
for relocated surgeries. Although the absolute number of relocated cases
is small, the higher proportion suggests that surgeries requiring a last
minute room change are more likely to run past the scheduled shift
boundary.

This difference likely reflects the circumstances under which
relocations occur. Cases that are reassigned to another room often face
tighter scheduling constraints, may start later than intended, or are
moved to resolve operational conflicts, all of which increase the chance
of extending into the evening.

Even so, the overall contribution of relocations to after hours workload
remains limited. Because relocations represent only a small fraction of
total surgical activity, their higher individual likelihood of overtime
does not translate into a large effect at the campus level.

\begin{longtable}[]{@{}
  >{\centering\arraybackslash}p{(\linewidth - 4\tabcolsep) * \real{0.2222}}
  >{\centering\arraybackslash}p{(\linewidth - 4\tabcolsep) * \real{0.1806}}
  >{\centering\arraybackslash}p{(\linewidth - 4\tabcolsep) * \real{0.4306}}@{}}
\caption{Share of after-hours surgeries by relocation status (Cathlab).
\label{tab:overtime_relocation}}\tabularnewline
\toprule\noalign{}
\begin{minipage}[b]{\linewidth}\centering
SwapGroup
\end{minipage} & \begin{minipage}[b]{\linewidth}\centering
AfterHours
\end{minipage} & \begin{minipage}[b]{\linewidth}\centering
Percentage of surgeries with afterhours
\end{minipage} \\
\midrule\noalign{}
\endfirsthead
\toprule\noalign{}
\begin{minipage}[b]{\linewidth}\centering
SwapGroup
\end{minipage} & \begin{minipage}[b]{\linewidth}\centering
AfterHours
\end{minipage} & \begin{minipage}[b]{\linewidth}\centering
Percentage of surgeries with afterhours
\end{minipage} \\
\midrule\noalign{}
\endhead
\bottomrule\noalign{}
\endlastfoot
Not relocated & 524 & 5.7 \\
Relocated & 12 & 10.6 \\
\end{longtable}

The distribution of overtime duration for relocated surgeries, shown in
Figure \ref{fig:overtime_relocated}, is concentrated at the lower end of
the range. Most relocated cases that run past the scheduled shift extend
by less than an hour, and several fall below the 30 minute mark. Only a
small number of surgeries exceed 60 minutes of overtime, and cases
beyond two hours are rare.

This pattern reinforces the earlier observation that relocated surgeries
tend to be short, flexible procedures. Once reassigned to another room,
they are usually completed without generating substantial after hours
workload. The sparse right tail indicates that although long overruns
can occur, they are exceptional rather than characteristic of relocated
cases.

Operationally, the distribution suggests that room swaps are handled
efficiently: they solve short term scheduling conflicts without creating
notable additional overtime.

\begin{figure}[H]

{\centering \includegraphics[width=0.7\linewidth,]{In-Depth_Analysis_Cathlab_files/figure-latex/overtime_relocated-1} 

}

\caption{Distribution of overtime duration for relocated surgeries \label{fig:overtime_relocated}}\label{fig:overtime_relocated}
\end{figure}

\newpage

\section{How do urgent cases disrupt elective
planning?}\label{how-do-urgent-cases-disrupt-elective-planning}

\subsection{When and where do urgent surgeries
occur?}\label{when-and-where-do-urgent-surgeries-occur}

The surgical caseload at Cathlab is overwhelmingly elective, as shown in
Table \ref{tab:urgency_dist}. Elective surgeries account for 97.5
percent of all recorded cases, while non elective procedures represent
only 2.5 percent. This distribution indicates that the operating
schedule is largely driven by planned activity, with urgent cases
forming a small but operationally important minority.

Even though non elective surgeries constitute a limited share, their
unpredictable timing and variable duration can influence daily workflow,
room allocation, and schedule stability. Understanding how these cases
are distributed across days, rooms, and time periods is essential for
assessing their impact on deviations, overtime, and planning robustness.

\begin{longtable}[]{@{}
  >{\centering\arraybackslash}p{(\linewidth - 4\tabcolsep) * \real{0.2222}}
  >{\centering\arraybackslash}p{(\linewidth - 4\tabcolsep) * \real{0.1111}}
  >{\centering\arraybackslash}p{(\linewidth - 4\tabcolsep) * \real{0.1528}}@{}}
\caption{Distribution of surgeries by urgency classification (Cathlab).
\label{tab:urgency_dist}}\tabularnewline
\toprule\noalign{}
\begin{minipage}[b]{\linewidth}\centering
UrgencyType
\end{minipage} & \begin{minipage}[b]{\linewidth}\centering
Total
\end{minipage} & \begin{minipage}[b]{\linewidth}\centering
SharePct
\end{minipage} \\
\midrule\noalign{}
\endfirsthead
\toprule\noalign{}
\begin{minipage}[b]{\linewidth}\centering
UrgencyType
\end{minipage} & \begin{minipage}[b]{\linewidth}\centering
Total
\end{minipage} & \begin{minipage}[b]{\linewidth}\centering
SharePct
\end{minipage} \\
\midrule\noalign{}
\endhead
\bottomrule\noalign{}
\endlastfoot
electief & 9050 & 97.5 \\
niet-electief & 232 & 2.5 \\
\end{longtable}

The distribution of non elective surgery start times, shown in Figure
\ref{fig:urgency_timing}, is clearly concentrated during daytime hours.
Activity begins to rise in the late morning and reaches its highest
density between roughly 13:00 and 15:00, indicating that many urgent
cases are handled during periods when elective schedules are already
active. After the afternoon peak, the density tapers gradually into the
early evening, reflecting the arrival of additional urgent cases later
in the day.

During night hours, the density falls to very low levels, with only
sporadic non elective surgeries starting between midnight and early
morning. Overall, the pattern shows that urgent procedures are not
uniformly distributed across the day but cluster around the core
operating period, where their overlap with elective activity has the
greatest potential to affect schedule stability and room availability.

\begin{figure}[H]

{\centering \includegraphics[width=0.7\linewidth,]{In-Depth_Analysis_Cathlab_files/figure-latex/urgency_timing-1} 

}

\caption{Daily distribution of non-elective surgery start times \label{fig:urgency_timing}}\label{fig:urgency_timing}
\end{figure}

The hourly distribution of non elective surgery start times follows a
remarkably consistent shape across weekdays, as illustrated in Figure
\ref{fig:urgency_weekday}. Each day shows a gradual rise in activity
during the late morning, followed by a pronounced peak in the early
afternoon, typically between 13:00 and 15:00. This clustering indicates
that urgent procedures are most often initiated during the same hours in
which elective activity is already at its busiest.

Although the overall volume differs by weekday, the temporal profile
remains similar: limited activity in the early morning, a sharp increase
toward midday, and a tapering into the early evening.

Weekend patterns are broadly comparable, though slightly more diffuse.
Both Saturday and Sunday display sustained activity throughout the day,
with an afternoon concentration still visible but less sharply defined
than during the week. The presence of urgent surgeries well into the
evening on weekends reflects the continuous nature of acute surgical
demand even in the absence of routine elective lists.

Taken together, the curves show that non elective cases rarely occur
overnight and instead cluster strongly within regular operating hours.
This alignment places additional pressure on shared capacity during peak
daytime periods, where both planned and unplanned surgical work
converge.

\begin{figure}[H]

{\centering \includegraphics[width=0.7\linewidth,]{In-Depth_Analysis_Cathlab_files/figure-latex/urgency_weekday-1} 

}

\caption{Hourly distribution of non-elective surgery start times by weekday \label{fig:urgency_weekday}}\label{fig:urgency_weekday}
\end{figure}

Elective surgeries dominate the workload on weekdays, as shown in Table
\ref{tab:urgency_wd}. From Monday through Friday, elective cases
consistently account for over 96 percent of all procedures, while non
elective surgeries make up only 1.6 to 3.1 percent. This narrow range
indicates a highly stable pattern of urgent demand during the regular
workweek, with very little day to day fluctuation. The weekday
environment is therefore overwhelmingly elective, and urgent procedures
represent only a small, predictable share of daily activity.

On the weekend, the case mix changes substantially. On Saturday,
elective surgeries fall to 79.7 percent, and non elective procedures
rise sharply to 20.3 percent. Sunday shows a similar pattern, with 76.7
percent elective and 23.3 percent non elective cases. These proportions
reflect the reduced volume of scheduled elective work combined with a
relatively steady flow of urgent cases that continue throughout the
week.

Overall, the data reveal a clear weekday weekend contrast: • Weekdays
are characterized by near total elective activity with minimal urgent
interruptions. • Weekends show a much higher share of urgent work, even
though absolute volumes are lower.

This shift underscores the need for flexible weekend staffing and
capacity, while weekday operations can largely focus on executing the
elective program with only minor urgent inflow.

\begin{longtable}[]{@{}
  >{\centering\arraybackslash}p{(\linewidth - 8\tabcolsep) * \real{0.1282}}
  >{\centering\arraybackslash}p{(\linewidth - 8\tabcolsep) * \real{0.1923}}
  >{\centering\arraybackslash}p{(\linewidth - 8\tabcolsep) * \real{0.2436}}
  >{\centering\arraybackslash}p{(\linewidth - 8\tabcolsep) * \real{0.1923}}
  >{\centering\arraybackslash}p{(\linewidth - 8\tabcolsep) * \real{0.2436}}@{}}
\caption{Weekday distribution of elective vs.~non-elective surgeries
(Cathlabs). \label{tab:urgency_wd}}\tabularnewline
\toprule\noalign{}
\begin{minipage}[b]{\linewidth}\centering
Weekday
\end{minipage} & \begin{minipage}[b]{\linewidth}\centering
Elective (n)
\end{minipage} & \begin{minipage}[b]{\linewidth}\centering
Non-elective (n)
\end{minipage} & \begin{minipage}[b]{\linewidth}\centering
Elective (\%)
\end{minipage} & \begin{minipage}[b]{\linewidth}\centering
Non-elective (\%)
\end{minipage} \\
\midrule\noalign{}
\endfirsthead
\toprule\noalign{}
\begin{minipage}[b]{\linewidth}\centering
Weekday
\end{minipage} & \begin{minipage}[b]{\linewidth}\centering
Elective (n)
\end{minipage} & \begin{minipage}[b]{\linewidth}\centering
Non-elective (n)
\end{minipage} & \begin{minipage}[b]{\linewidth}\centering
Elective (\%)
\end{minipage} & \begin{minipage}[b]{\linewidth}\centering
Non-elective (\%)
\end{minipage} \\
\midrule\noalign{}
\endhead
\bottomrule\noalign{}
\endlastfoot
Ma & 1355 & 43 & 96.9 & 3.1 \\
Di & 1777 & 35 & 98.1 & 1.9 \\
Wo & 2261 & 47 & 98 & 2 \\
Do & 2001 & 33 & 98.4 & 1.6 \\
Vr & 1553 & 45 & 97.2 & 2.8 \\
Za & 47 & 12 & 79.7 & 20.3 \\
Zo & 56 & 17 & 76.7 & 23.3 \\
\end{longtable}

\subsection{Do urgent cases drive
overtime?}\label{do-urgent-cases-drive-overtime}

Elective and non elective surgeries differ substantially in how often
they extend beyond scheduled shift hours, as shown in Table
\ref{tab:afterhours_urgency}. Only 5.5 percent of elective surgeries run
after hours, whereas 19.4 percent of non elective procedures do so.
Although urgent cases make up a small fraction of total activity, they
are three times more likely to continue beyond 16:30, confirming their
disproportionate contribution to after hours workload.

The duration of overtime also differs sharply. Elective procedures that
run past shift boundaries do so for a median of 2.2 minutes, with a 95th
percentile of 5 minutes, indicating that most elective overruns are
short and tightly controlled. In contrast, non elective surgeries have a
mean overtime of 10.3 minutes and a 95th percentile of 77.3 minutes,
showing that when urgent cases extend beyond the shift, they may require
significantly more time.

Together, these results highlight a dual effect: • Elective surgeries
generate most after hours cases in absolute numbers, simply because they
are far more frequent. • Non elective surgeries generate more intense
after hours demand, being both more likely to run late and more likely
to do so for extended periods.

The findings underline the operational importance of maintaining
flexible evening capacity to accommodate the unpredictability inherent
in urgent surgical care.

\begin{longtable}[]{@{}
  >{\centering\arraybackslash}p{(\linewidth - 10\tabcolsep) * \real{0.1376}}
  >{\centering\arraybackslash}p{(\linewidth - 10\tabcolsep) * \real{0.1101}}
  >{\centering\arraybackslash}p{(\linewidth - 10\tabcolsep) * \real{0.1651}}
  >{\centering\arraybackslash}p{(\linewidth - 10\tabcolsep) * \real{0.1651}}
  >{\centering\arraybackslash}p{(\linewidth - 10\tabcolsep) * \real{0.2018}}
  >{\centering\arraybackslash}p{(\linewidth - 10\tabcolsep) * \real{0.2202}}@{}}
\caption{After-hours frequency and overtime duration by urgency type
(Cathlab). \label{tab:afterhours_urgency}}\tabularnewline
\toprule\noalign{}
\begin{minipage}[b]{\linewidth}\centering
Urgency type
\end{minipage} & \begin{minipage}[b]{\linewidth}\centering
Total (n)
\end{minipage} & \begin{minipage}[b]{\linewidth}\centering
After-hours (n)
\end{minipage} & \begin{minipage}[b]{\linewidth}\centering
After-hours (\%)
\end{minipage} & \begin{minipage}[b]{\linewidth}\centering
Mean overtime (min)
\end{minipage} & \begin{minipage}[b]{\linewidth}\centering
95th percentile (min)
\end{minipage} \\
\midrule\noalign{}
\endfirsthead
\toprule\noalign{}
\begin{minipage}[b]{\linewidth}\centering
Urgency type
\end{minipage} & \begin{minipage}[b]{\linewidth}\centering
Total (n)
\end{minipage} & \begin{minipage}[b]{\linewidth}\centering
After-hours (n)
\end{minipage} & \begin{minipage}[b]{\linewidth}\centering
After-hours (\%)
\end{minipage} & \begin{minipage}[b]{\linewidth}\centering
Mean overtime (min)
\end{minipage} & \begin{minipage}[b]{\linewidth}\centering
95th percentile (min)
\end{minipage} \\
\midrule\noalign{}
\endhead
\bottomrule\noalign{}
\endlastfoot
Elective & 9050 & 494 & 5.5 & 2.2 & 5 \\
NonElective & 232 & 45 & 19.4 & 10.3 & 77.3 \\
\end{longtable}

\subsection{Do urgent cases interrupt elective
planning?}\label{do-urgent-cases-interrupt-elective-planning}

The results in Table \ref{tab:overlap_swap} show that room swaps among
elective surgeries are extremely rare and occur almost exclusively in
cases without an overlap. Among the 8,944 elective surgeries with no
overlap, only 105 cases involved a room swap, corresponding to 1.2
percent. In contrast, none of the 100 elective surgeries that overlapped
with a non elective case were relocated.

This pattern suggests two important points. First, room reallocations in
the elective program are uncommon and appear only in isolated
situations. Second, when an elective case overlaps with an urgent
intervention, the system absorbs this pressure without resorting to
moving the elective case to another operating room. Overlaps are
therefore handled through local adjustments in timing, sequencing, or
resource coordination rather than through physical room changes.

Overall, the data indicate that urgent activity does not materially
disrupt room assignments within the elective workflow. The room
allocation process remains stable even when unplanned demand intersects
with scheduled activity.

\begin{longtable}[]{@{}
  >{\centering\arraybackslash}p{(\linewidth - 6\tabcolsep) * \real{0.1806}}
  >{\centering\arraybackslash}p{(\linewidth - 6\tabcolsep) * \real{0.0972}}
  >{\centering\arraybackslash}p{(\linewidth - 6\tabcolsep) * \real{0.1667}}
  >{\centering\arraybackslash}p{(\linewidth - 6\tabcolsep) * \real{0.2083}}@{}}
\caption{Room swap frequency among elective cases: overlapped vs.~not
overlapped. \label{tab:overlap_swap}}\tabularnewline
\toprule\noalign{}
\begin{minipage}[b]{\linewidth}\centering
Overlapped
\end{minipage} & \begin{minipage}[b]{\linewidth}\centering
N
\end{minipage} & \begin{minipage}[b]{\linewidth}\centering
RoomSwaps
\end{minipage} & \begin{minipage}[b]{\linewidth}\centering
ShareSwapped
\end{minipage} \\
\midrule\noalign{}
\endfirsthead
\toprule\noalign{}
\begin{minipage}[b]{\linewidth}\centering
Overlapped
\end{minipage} & \begin{minipage}[b]{\linewidth}\centering
N
\end{minipage} & \begin{minipage}[b]{\linewidth}\centering
RoomSwaps
\end{minipage} & \begin{minipage}[b]{\linewidth}\centering
ShareSwapped
\end{minipage} \\
\midrule\noalign{}
\endhead
\bottomrule\noalign{}
\endlastfoot
No & 8944 & 105 & 1.2 \\
Yes & 100 & 0 & 0 \\
\end{longtable}

As shown in Table \ref{tab:overlap_days}, overlaps between urgent and
elective surgeries are relatively uncommon at the day level. Out of 963
observed operating days, 46 days contained at least one instance where a
non elective case overlapped with an elective case in the same operating
room. This corresponds to 4.8 percent of all days.

This low percentage indicates that while urgent cases do occasionally
intersect with scheduled elective activity, such overlaps occur on only
a small fraction of days. Rather than being a routine operational
pattern, they represent isolated events that arise intermittently. When
overlaps do occur, they may require real time coordination, but their
infrequency suggests that they are not a dominant or persistent source
of scheduling disruption in the daily workflow.

Overall, the data show that urgent elective overlaps are operationally
relevant but not a regular or pervasive feature of day to day activity.

\begin{longtable}[]{@{}
  >{\centering\arraybackslash}p{(\linewidth - 4\tabcolsep) * \real{0.1944}}
  >{\centering\arraybackslash}p{(\linewidth - 4\tabcolsep) * \real{0.1667}}
  >{\centering\arraybackslash}p{(\linewidth - 4\tabcolsep) * \real{0.1667}}@{}}
\caption{Days with at least one urgent--elective overlap in the same OR.
\label{tab:overlap_days}}\tabularnewline
\toprule\noalign{}
\begin{minipage}[b]{\linewidth}\centering
OverlapDays
\end{minipage} & \begin{minipage}[b]{\linewidth}\centering
TotalDays
\end{minipage} & \begin{minipage}[b]{\linewidth}\centering
SharePct
\end{minipage} \\
\midrule\noalign{}
\endfirsthead
\toprule\noalign{}
\begin{minipage}[b]{\linewidth}\centering
OverlapDays
\end{minipage} & \begin{minipage}[b]{\linewidth}\centering
TotalDays
\end{minipage} & \begin{minipage}[b]{\linewidth}\centering
SharePct
\end{minipage} \\
\midrule\noalign{}
\endhead
\bottomrule\noalign{}
\endlastfoot
46 & 963 & 4.8 \\
\end{longtable}

The proportion of elective surgeries affected by urgent overlaps varies
from month to month but remains very low overall, typically between 0.5
and 2.5 percent, as shown in Figure \ref{fig:overlap_trend}. Only a few
isolated months reach values near 3 percent. Rather than forming a clear
upward or downward trend, the series oscillates irregularly, with
short-lived peaks followed by returns to lower levels.

This pattern suggests that urgent overlaps are small in scale and
largely driven by incidental fluctuations in daily operational
conditions. The absence of a persistent increase over time indicates
that the interaction between urgent and elective activity is stable, and
overlaps occur only in a small fraction of elective cases each month.

\begin{figure}[H]

{\centering \includegraphics[width=0.7\linewidth,]{In-Depth_Analysis_Cathlab_files/figure-latex/overlap_trend-1} 

}

\caption{Monthly share of elective surgeries affected by urgent overlaps \label{fig:overlap_trend}}\label{fig:overlap_trend}
\end{figure}

As shown in Figure \ref{fig:overlap_hourly}, urgent--elective overlaps
are concentrated almost entirely within standard daytime operating
hours. The first overlaps appear shortly after 08:00 and increase
steadily through the morning. The frequency reaches its highest levels
between 13:00 and 15:00, where the histogram shows the most pronounced
bars, indicating that these afternoon hours represent the peak period
for schedule conflicts. After 16:00 the number of overlaps falls
sharply, and almost no events occur outside the regular daytime window.

This temporal concentration closely mirrors the daily timing of
non-elective surgery starts, with both peaking during early and
mid-afternoon. The alignment of these patterns suggests that urgent
cases most often arise at the exact times when elective activity is
already at full capacity, thereby heightening the likelihood of
operational interference. The clustering around midday underscores how
overlaps are fundamentally a product of simultaneous demand during the
busiest phase of the surgical schedule rather than a result of
late-evening or night-time emergencies.

\begin{figure}[H]

{\centering \includegraphics[width=0.7\linewidth,]{In-Depth_Analysis_Cathlab_files/figure-latex/overlap_hourly-1} 

}

\caption{Hourly frequency of urgent–elective overlaps \label{fig:overlap_hourly}}\label{fig:overlap_hourly}
\end{figure}

Figure \ref{fig:overlap_delay} shows that elective surgeries affected by
an urgent overlap experience far greater variability in their start
times than those without an overlap. The non-overlapped group remains
stable across the entire period, with average start deviations
fluctuating only a few minutes around zero. In contrast, the overlapped
group displays large swings from month to month. Several months show
substantial negative deviations, hundreds of minutes earlier than
planned, driven by a small number of extreme cases. At other points,
overlapped surgeries start slightly later than expected, but the pattern
is irregular rather than consistently delayed.

These intermittent spikes suggest that overlaps do not uniformly shift
elective cases later. Instead, they create volatile and occasionally
extreme deviations, reflecting reactive adjustments made in response to
urgent demand. The stability of the non-overlapped group highlights that
these outliers are tightly linked to overlap events rather than to
broader scheduling behavior. Overall, the figure demonstrates that
urgent overlaps introduce unpredictability into elective start times,
not through systematic delay but through sporadic and substantial
departures from the planned schedule.

\begin{figure}[H]

{\centering \includegraphics[width=0.7\linewidth,]{In-Depth_Analysis_Cathlab_files/figure-latex/overlap_delay-1} 

}

\caption{Average start delay of elective surgeries by overlap status \label{fig:overlap_delay}}\label{fig:overlap_delay}
\end{figure}

Table \ref{tab:or_overlap} shows that urgent--elective overlaps are
concentrated in only a few operating rooms. KO05 accounts for the
largest number of overlap events (71 cases), which is 2.7 percent of all
elective surgeries performed in that room. KO06 follows with 15 overlaps
(1.8 percent), while KO02, KO07 and KO01 each record only a very small
number of affected cases. The pattern indicates that overlaps are not
evenly distributed across the OR complex but occur primarily in rooms
with higher elective throughput or mixed case profiles.

The median start-time effects also vary substantially by room. In KO05
and KO01, overlapped elective surgeries start earlier on average than
non-overlapped ones (−14 and −127 minutes, respectively), driven by a
few sizeable early starts. In KO02 the median shift is also negative
(−64 minutes). By contrast, KO07 shows a modest positive delay of 22.5
minutes, and KO06 exhibits a median delay of −11 minutes despite a large
positive reference value for non-overlapped cases (90 minutes). These
heterogeneous outcomes suggest that the operational consequences of
overlaps depend strongly on local workflow dynamics rather than
following a uniform pattern.

Taken together, the results show that overlaps remain infrequent across
most rooms but can produce substantial deviations---either earlier or
later---when they occur. The concentration of overlap events in a small
subset of ORs highlights the importance of room-specific workflow
characteristics in shaping how urgent cases interact with the elective
schedule.

\begin{longtable}[]{@{}
  >{\centering\arraybackslash}p{(\linewidth - 10\tabcolsep) * \real{0.1165}}
  >{\centering\arraybackslash}p{(\linewidth - 10\tabcolsep) * \real{0.0777}}
  >{\centering\arraybackslash}p{(\linewidth - 10\tabcolsep) * \real{0.1262}}
  >{\centering\arraybackslash}p{(\linewidth - 10\tabcolsep) * \real{0.1262}}
  >{\centering\arraybackslash}p{(\linewidth - 10\tabcolsep) * \real{0.2913}}
  >{\centering\arraybackslash}p{(\linewidth - 10\tabcolsep) * \real{0.2621}}@{}}
\caption{Operating rooms with urgent--elective overlap.
\label{tab:or_overlap}}\tabularnewline
\toprule\noalign{}
\begin{minipage}[b]{\linewidth}\centering
PlannedOR
\end{minipage} & \begin{minipage}[b]{\linewidth}\centering
Count
\end{minipage} & \begin{minipage}[b]{\linewidth}\centering
n\_elective
\end{minipage} & \begin{minipage}[b]{\linewidth}\centering
Percentage
\end{minipage} & \begin{minipage}[b]{\linewidth}\centering
Med StartDelay (No Overlap)
\end{minipage} & \begin{minipage}[b]{\linewidth}\centering
Med StartDelay (Overlap)
\end{minipage} \\
\midrule\noalign{}
\endfirsthead
\toprule\noalign{}
\begin{minipage}[b]{\linewidth}\centering
PlannedOR
\end{minipage} & \begin{minipage}[b]{\linewidth}\centering
Count
\end{minipage} & \begin{minipage}[b]{\linewidth}\centering
n\_elective
\end{minipage} & \begin{minipage}[b]{\linewidth}\centering
Percentage
\end{minipage} & \begin{minipage}[b]{\linewidth}\centering
Med StartDelay (No Overlap)
\end{minipage} & \begin{minipage}[b]{\linewidth}\centering
Med StartDelay (Overlap)
\end{minipage} \\
\midrule\noalign{}
\endhead
\bottomrule\noalign{}
\endlastfoot
KO05 & 71 & 2597 & 2.7 & 1 & -14 \\
KO06 & 15 & 818 & 1.8 & 90 & -11 \\
KO02 & 7 & 2835 & 0.2 & 4 & -64 \\
KO07 & 4 & 896 & 0.4 & 13 & 22.5 \\
KO01 & 3 & 1854 & 0.2 & 2 & -127 \\
\end{longtable}

\newpage

\section{How stable is the surgical schedule in
execution?}\label{how-stable-is-the-surgical-schedule-in-execution}

Beyond accurate planning, the stability of the realised schedule is a
key indicator of operational performance in the operating theatre. Even
when planned start times, durations, and room assignments are correct,
day-to-day execution may deviate because of unforeseen delays, emergency
insertions, or coordination issues.

This section evaluates how consistently the planned schedule is
maintained throughout the day by focusing on two complementary aspects:

\begin{itemize}
\tightlist
\item
  \emph{Shift changes}, where surgeries are executed in a different
  shift than originally planned, indicating how often the daily
  programme must be adjusted in real time.
\item
  \emph{Idle or gap times}, representing the pauses between consecutive
  surgeries within the same room, which reflect how efficiently the
  available capacity is utilised.
\end{itemize}

Together, these indicators capture the operational resilience of the OR
complex and the ability to absorb disruptions and maintain smooth
throughput despite variable circumstances. High rates of shift changes
or long idle periods point to latent inefficiencies in daily planning
and coordination, whereas stability across these metrics signals a
well-synchronised workflow.

\subsection{How often are surgeries moved to another
shift?}\label{how-often-are-surgeries-moved-to-another-shift}

Shift changes occur when a surgery is performed in a different operating
shift than originally planned. This happens when cases start much later
than scheduled or when they are postponed into a later block of the day.
Actual shift assignment is determined using the observed start time
relative to the standard shift boundaries at Cathlab: day
(08:00--16:30), evening (16:30--22:00), and night (22:00--08:00).

Table \ref{tab:shift_mismatch} shows that 325 surgeries, about 3.5
percent of all recorded cases, were ultimately performed in a different
shift than planned. These mismatches are associated with substantial
start delays, averaging 169.7 minutes, indicating that a shift change is
rarely the result of a minor deviation but rather of a major
postponement within the daily program.

Interestingly, once these surgeries do begin, their observed duration
tends to be shorter than estimated: the mean duration difference is
--87.6 minutes. This suggests that shift changes are often linked to
scheduling or sequencing issues rather than to unexpectedly long
procedures.

Average overtime among these cases is 11 minutes, showing that most
shift mismatches are driven by late starts rather than by large overruns
extending deep into the next shift.

Although these events represent a small fraction of total surgical
volume, their large delays and shift transitions can create
disproportionate downstream effects on staffing, turnover, and room
availability, making them an important operational phenomenon to
monitor.

\begin{longtable}[]{@{}
  >{\centering\arraybackslash}p{(\linewidth - 10\tabcolsep) * \real{0.0609}}
  >{\centering\arraybackslash}p{(\linewidth - 10\tabcolsep) * \real{0.1391}}
  >{\centering\arraybackslash}p{(\linewidth - 10\tabcolsep) * \real{0.1565}}
  >{\centering\arraybackslash}p{(\linewidth - 10\tabcolsep) * \real{0.2174}}
  >{\centering\arraybackslash}p{(\linewidth - 10\tabcolsep) * \real{0.2348}}
  >{\centering\arraybackslash}p{(\linewidth - 10\tabcolsep) * \real{0.1913}}@{}}
\caption{Cases performed in a different shift than planned (Cathlab).
\label{tab:shift_mismatch}}\tabularnewline
\toprule\noalign{}
\begin{minipage}[b]{\linewidth}\centering
n
\end{minipage} & \begin{minipage}[b]{\linewidth}\centering
MismatchCount
\end{minipage} & \begin{minipage}[b]{\linewidth}\centering
Share moved (\%)
\end{minipage} & \begin{minipage}[b]{\linewidth}\centering
Mean Start Delay (min)
\end{minipage} & \begin{minipage}[b]{\linewidth}\centering
Mean Duration Diff (min)
\end{minipage} & \begin{minipage}[b]{\linewidth}\centering
Mean Overtime (min)
\end{minipage} \\
\midrule\noalign{}
\endfirsthead
\toprule\noalign{}
\begin{minipage}[b]{\linewidth}\centering
n
\end{minipage} & \begin{minipage}[b]{\linewidth}\centering
MismatchCount
\end{minipage} & \begin{minipage}[b]{\linewidth}\centering
Share moved (\%)
\end{minipage} & \begin{minipage}[b]{\linewidth}\centering
Mean Start Delay (min)
\end{minipage} & \begin{minipage}[b]{\linewidth}\centering
Mean Duration Diff (min)
\end{minipage} & \begin{minipage}[b]{\linewidth}\centering
Mean Overtime (min)
\end{minipage} \\
\midrule\noalign{}
\endhead
\bottomrule\noalign{}
\endlastfoot
9282 & 325 & 3.5 & 169.7 & -87.6 & 11 \\
\end{longtable}

The monthly share of surgeries performed in a different shift shows
noticeable short-term variability but no structural worsening over time,
as illustrated in Figure \ref{fig:shift_moved_monthly}. In the first
months of 2022, the proportion fluctuates around 4 to 6 percent,
representing the highest levels observed in the series. After this
initial period, the share settles into a more stable range between 3 and
5 percent, with only occasional brief peaks.

Across 2023 and 2024 the monthly values remain within this narrow band,
and the pattern in early 2025 is consistent with the previous years.
Although the line oscillates from month to month, the movements appear
cyclical rather than directional. There is no indication of a sustained
upward trend or structural deterioration in shift reassignments.
Instead, the variation likely reflects temporary scheduling pressures or
day-to-day fluctuations in case timing rather than long-term changes in
planning performance.

\begin{figure}[H]

{\centering \includegraphics[width=0.7\linewidth,]{In-Depth_Analysis_Cathlab_files/figure-latex/shift_moved_monthly-1} 

}

\caption{Monthly percentage of surgeries moved to another shift \label{fig:shift_moved_monthly}}\label{fig:shift_moved_monthly}
\end{figure}

Overall, the data suggest a moderate improvement in schedule adherence
since 2022, followed by a period of relative stability. The
month-to-month variation likely reflects normal fluctuations in workload
and case mix rather than structural scheduling issues.

\subsection{How much idle time occurs between consecutive
surgeries?}\label{how-much-idle-time-occurs-between-consecutive-surgeries}

Idle or gap time represents the short interval between the end of one
surgery and the start of the next within the same operating room during
the regular day shift (08:00--16:30). It reflects how efficiently rooms
are turned over and how well consecutive cases are synchronised in
practice. For each room-day combination, gap time covers either the
delay between shift start and the first case or the transition between
consecutive surgeries, while pauses longer than sixty minutes are
excluded because they represent planned breaks rather than operational
idle time.

At Campus Cathlab, the average idle time between consecutive surgeries
is 7.1 minutes, with a median of 4 minutes, as shown in Table
\ref{tab:idle_summary}. This low central tendency indicates that
turnovers are generally fast. The interquartile range of 7 minutes shows
that most transitions occur within a tight window, with limited
variability across cases. Longer gaps are uncommon: the 95th percentile
is 23 minutes, and the maximum observed idle time is 60 minutes, the
threshold applied for excluding planned downtime.

Together, these figures demonstrate a consistently efficient turnover
process in the operating rooms, where most transitions are completed
within only a few minutes and extended idle periods occur only rarely.

\begin{longtable}[]{@{}
  >{\centering\arraybackslash}p{(\linewidth - 16\tabcolsep) * \real{0.0972}}
  >{\centering\arraybackslash}p{(\linewidth - 16\tabcolsep) * \real{0.1250}}
  >{\centering\arraybackslash}p{(\linewidth - 16\tabcolsep) * \real{0.0833}}
  >{\centering\arraybackslash}p{(\linewidth - 16\tabcolsep) * \real{0.0833}}
  >{\centering\arraybackslash}p{(\linewidth - 16\tabcolsep) * \real{0.0833}}
  >{\centering\arraybackslash}p{(\linewidth - 16\tabcolsep) * \real{0.0833}}
  >{\centering\arraybackslash}p{(\linewidth - 16\tabcolsep) * \real{0.0833}}
  >{\centering\arraybackslash}p{(\linewidth - 16\tabcolsep) * \real{0.0833}}
  >{\centering\arraybackslash}p{(\linewidth - 16\tabcolsep) * \real{0.1111}}@{}}
\caption{Summary statistics of idle time between consecutive surgeries
(minutes). \label{tab:idle_summary}}\tabularnewline
\toprule\noalign{}
\begin{minipage}[b]{\linewidth}\centering
Mean
\end{minipage} & \begin{minipage}[b]{\linewidth}\centering
Median
\end{minipage} & \begin{minipage}[b]{\linewidth}\centering
SD
\end{minipage} & \begin{minipage}[b]{\linewidth}\centering
IQR
\end{minipage} & \begin{minipage}[b]{\linewidth}\centering
Min
\end{minipage} & \begin{minipage}[b]{\linewidth}\centering
Max
\end{minipage} & \begin{minipage}[b]{\linewidth}\centering
P95
\end{minipage} & \begin{minipage}[b]{\linewidth}\centering
P99
\end{minipage} & \begin{minipage}[b]{\linewidth}\centering
P99.9
\end{minipage} \\
\midrule\noalign{}
\endfirsthead
\toprule\noalign{}
\begin{minipage}[b]{\linewidth}\centering
Mean
\end{minipage} & \begin{minipage}[b]{\linewidth}\centering
Median
\end{minipage} & \begin{minipage}[b]{\linewidth}\centering
SD
\end{minipage} & \begin{minipage}[b]{\linewidth}\centering
IQR
\end{minipage} & \begin{minipage}[b]{\linewidth}\centering
Min
\end{minipage} & \begin{minipage}[b]{\linewidth}\centering
Max
\end{minipage} & \begin{minipage}[b]{\linewidth}\centering
P95
\end{minipage} & \begin{minipage}[b]{\linewidth}\centering
P99
\end{minipage} & \begin{minipage}[b]{\linewidth}\centering
P99.9
\end{minipage} \\
\midrule\noalign{}
\endhead
\bottomrule\noalign{}
\endlastfoot
7.1 & 4 & 8.4 & 7 & 0 & 60 & 23 & 45 & 59 \\
\end{longtable}

As shown in Figure \ref{fig:idle_distribution}, the distribution
confirms that most idle intervals are extremely short. The largest bar
corresponds to gaps of less than five minutes, which represent a
substantial share of all transitions. When expanding to the first ten
minutes, the histogram shows that the majority of turnovers fall within
this range. Beyond that point, the frequency drops quickly: only a small
fraction of intervals exceed twenty minutes, and transitions longer than
thirty minutes become rare. The right tail remains thin, as pauses
beyond sixty minutes are excluded from the analysis.

The overall pattern illustrates a highly efficient turnover process.
Short idle periods dominate routine workflow, while longer gaps appear
as isolated events rather than evidence of systematic underutilisation.
Such a distribution is characteristic of a well-synchronised operating
environment in which room preparation and case sequencing proceed with
limited and predictable downtime.

\begin{figure}[H]

{\centering \includegraphics[width=0.7\linewidth,]{In-Depth_Analysis_Cathlab_files/figure-latex/idle_distribution-1} 

}

\caption{Distribution of idle time between surgeries (grouped per 5 min) \label{fig:idle_distribution}}\label{fig:idle_distribution}
\end{figure}

\subsection{How does idle time vary across operating
rooms?}\label{how-does-idle-time-vary-across-operating-rooms}

As shown in Figure \ref{fig:idle_or}, mean and median idle times vary
moderately across the six operating rooms included in the analysis. In
every room, the median idle time is noticeably shorter than the mean,
confirming that most turnover intervals are brief while a smaller number
of longer pauses increase the average. This consistent gap reflects the
right-skewed nature of idle-time distributions: short intervals dominate
daily workflow, but extended gaps still occur occasionally.

Absolute differences between rooms remain limited. Median idle times
cluster tightly around a few minutes, suggesting that turnover
performance is broadly comparable. Rooms with higher mean values, such
as KO07 and KO05, also show slightly elevated medians, indicating the
presence of more frequent or somewhat longer pauses. By contrast, rooms
like KO01 and KO04 record the shortest turnover intervals, pointing to
particularly streamlined case transitions.

Overall, the figure highlights a stable and efficient turnover process
across the OR complex. The modest variation between rooms suggests that
differences arise mainly from isolated longer gaps---often related to
case mix or operational sequencing---rather than from systematic
inefficiencies in room coordination.

\begin{figure}[H]

{\centering \includegraphics[width=0.7\linewidth,]{In-Depth_Analysis_Cathlab_files/figure-latex/idle_or-1} 

}

\caption{Mean and median idle time per operating room \label{fig:idle_or}}\label{fig:idle_or}
\end{figure}

As shown in Table \ref{tab:idle_variability}, the standard deviation
(SD) and coefficient of variation (CV) provide a detailed view of how
strongly idle times fluctuate across operating rooms. Most rooms display
SD values between 5 and 12 minutes, indicating a generally narrow spread
around their typical turnover interval. The associated CV values, which
range from about 83 to 136 percent for the main rooms, confirm that
although turnover times are short, they remain noticeably variable
relative to their mean.

Rooms such as KO03 and KO04 stand out with substantially higher CV
values, reaching 136 and 204 percent respectively. In both cases, the
elevated variability is linked to very small sample sizes and the
presence of a few widely scattered idle intervals. KO06 also shows a
relatively large SD (15 minutes), suggesting a broader distribution of
turnover times compared with the other high-volume rooms.

By contrast, the busiest rooms---KO05, KO02, and KO01---exhibit the
lowest SDs (5.5 to 7 minutes) and correspondingly lower CVs (83 to 129
percent). These patterns indicate that in high-throughput rooms,
turnover times tend to cluster more tightly around their typical values,
reflecting more consistent workflows and a larger number of
observations.

Overall, idle-time variability differs meaningfully across rooms but
remains contained for most of the OR complex. Short turnover periods
clearly dominate, while a few rooms---mostly those with limited case
counts---display broader distributions that reflect either irregular
workflows or small-sample volatility rather than structural
inefficiency.

\begin{longtable}[]{@{}
  >{\centering\arraybackslash}p{(\linewidth - 6\tabcolsep) * \real{0.1528}}
  >{\centering\arraybackslash}p{(\linewidth - 6\tabcolsep) * \real{0.1111}}
  >{\centering\arraybackslash}p{(\linewidth - 6\tabcolsep) * \real{0.2917}}
  >{\centering\arraybackslash}p{(\linewidth - 6\tabcolsep) * \real{0.3750}}@{}}
\caption{Standard deviation and coefficient of variation of idle time
per operating room (minutes), sorted by median idle time.
\label{tab:idle_variability}}\tabularnewline
\toprule\noalign{}
\begin{minipage}[b]{\linewidth}\centering
ActualOR
\end{minipage} & \begin{minipage}[b]{\linewidth}\centering
Cases
\end{minipage} & \begin{minipage}[b]{\linewidth}\centering
Standard deviation
\end{minipage} & \begin{minipage}[b]{\linewidth}\centering
Coefficient of variation
\end{minipage} \\
\midrule\noalign{}
\endfirsthead
\toprule\noalign{}
\begin{minipage}[b]{\linewidth}\centering
ActualOR
\end{minipage} & \begin{minipage}[b]{\linewidth}\centering
Cases
\end{minipage} & \begin{minipage}[b]{\linewidth}\centering
Standard deviation
\end{minipage} & \begin{minipage}[b]{\linewidth}\centering
Coefficient of variation
\end{minipage} \\
\midrule\noalign{}
\endhead
\bottomrule\noalign{}
\endlastfoot
KO03 & 24 & 17.5 & 136 \\
KO06 & 506 & 15 & 120.2 \\
KO07 & 571 & 11.4 & 104.1 \\
KO04 & 9 & 19.1 & 204.1 \\
KO05 & 2064 & 7 & 83.4 \\
KO02 & 2216 & 6.5 & 129.8 \\
KO01 & 1442 & 5.5 & 113.2 \\
\end{longtable}

\section{Conclusion}\label{conclusion}

This analysis of Campus Cathlab provides a detailed and data-driven
account of how surgical activity unfolds from planning to execution
within the operating theatre. By examining timing accuracy, room
allocation, after-hours activity, urgent case dynamics, shift
transitions, and turnover performance, the report highlights both the
system's strengths and its operational sensitivities.

Across sections, planned and observed durations align well on average,
though notable variability persists between procedure types. Start-time
delays are common and tend to accumulate throughout the day, with 58.3
percent of surgeries starting later than planned, with a mean delay of
69.6 minutes and a median delay of 28 minutes. These delays can push
finishing times later than scheduled, with extreme delays reaching up to
1,437 minutes. While the overall program remains reliable, these small
deviations illustrate how tightly sequenced lists react to routine
sources of disruption.

After-hours activity forms a recurring but contained part of the
workload, with around 5.8 percent of surgeries extending beyond regular
operating hours. Most overruns are modest, with the mean overtime
duration being 41.8 minutes and the median overtime being 32 minutes.
However, a small number of complex procedures account for a
disproportionate share of total overtime, with some cases extending up
to 108.3 minutes (P95). Room assignments are highly stable, with roughly
1.2 percent of surgeries being relocated to a different operating room,
typically short and low-complexity cases that impose minimal downstream
effects.

Urgent cases represent about 2.5 percent of activity but frequently
coincide with elective work. Although their share is modest, overlaps
occur on nearly 70 percent of days, underscoring their continuous
operational relevance. These overlaps can delay elective starts but
rarely trigger broader rescheduling or widespread room reallocation,
indicating effective day-to-day coordination between planned and
unplanned care.

Shift transitions and idle times point to an efficient process. About
3.5 percent of surgeries occur in a different shift than originally
planned, a small but operationally meaningful category linked to
substantial delays earlier in the day, with a mean start delay of 169.7
minutes. Idle periods between cases remain short and consistent across
rooms, with a median of 4 minutes and a mean of 7.1 minutes. Turnover is
therefore fast and well-stabilized, although a few rooms show broader
distributions reflecting more diverse local workflows.

Taken together, the findings present a comprehensive picture of
operating room performance at Campus Cathlab. Surgical planning
generally translates well into real-world execution, with a system that
balances predictability and flexibility in a high-throughput
environment. Opportunities for improvement remain modest and
concentrated: stabilizing start times, containing small day-level
delays, and maintaining robust coordination during periods of urgent
demand.

\end{document}
